\documentclass[11pt,a4paper]{book}
\usepackage{libertine}
\usepackage{fontspec}
\IfFontExistsTF{Menlo}{%
  \setmonofont{Menlo}[Scale=MatchLowercase]%
}{%
  \IfFontExistsTF{Latin Modern Mono}{%
    \setmonofont{Latin Modern Mono}[Scale=MatchLowercase]%
  }{%
    \IfFontExistsTF{Courier New}{%
      \setmonofont{Courier New}[Scale=MatchLowercase]%
    }{}%
  }%
}
\usepackage{newunicodechar}
\usepackage{booktabs}

% Map Agda symbols to math mode
\newunicodechar{∀}{\ensuremath{\forall}}
\newunicodechar{ℏ}{\ensuremath{\hbar}}
\newunicodechar{→}{\ensuremath{\to}}
\newunicodechar{≡}{\ensuremath{\equiv}}
\newunicodechar{ℕ}{\ensuremath{\mathbb{N}}}
\newunicodechar{ℤ}{\ensuremath{\mathbb{Z}}}
\newunicodechar{ℚ}{\ensuremath{\mathbb{Q}}}
\newunicodechar{∸}{\ensuremath{\dot{-}}}
\newunicodechar{≤}{\ensuremath{\le}}
\newunicodechar{≥}{\ensuremath{\ge}}
\newunicodechar{∘}{\ensuremath{\circ}}
\newunicodechar{×}{\ensuremath{\times}}
\newunicodechar{⊎}{\ensuremath{\uplus}}
\newunicodechar{⊥}{\ensuremath{\bot}}
\newunicodechar{⊤}{\ensuremath{\top}}
\newunicodechar{¬}{\ensuremath{\neg}}
\newunicodechar{≢}{\ensuremath{\not\equiv}}
\newunicodechar{○}{\ensuremath{\circ}}
\newunicodechar{●}{\ensuremath{\bullet}}
\newunicodechar{∷}{\ensuremath{::}}
\newunicodechar{λ}{\ensuremath{\lambda}}
\newunicodechar{₀}{\ensuremath{_0}}
\newunicodechar{₁}{\ensuremath{_1}}
\newunicodechar{₂}{\ensuremath{_2}}
\newunicodechar{₃}{\ensuremath{_3}}
\newunicodechar{₄}{\ensuremath{_4}}
\newunicodechar{₅}{\ensuremath{_5}}
\newunicodechar{₆}{\ensuremath{_6}}
\newunicodechar{₇}{\ensuremath{_7}}
\newunicodechar{₈}{\ensuremath{_8}}
\newunicodechar{₉}{\ensuremath{_9}}
\newunicodechar{⊨}{\ensuremath{\models}}
\newunicodechar{≈}{\ensuremath{\approx}}
\newunicodechar{≃}{\ensuremath{\simeq}}
\newunicodechar{≅}{\ensuremath{\cong}}
\newunicodechar{≟}{\ensuremath{\stackrel{?}{=}}}
\newunicodechar{′}{\ensuremath{'}}
\newunicodechar{″}{\ensuremath{''}}
\newunicodechar{‴}{\ensuremath{'''}}
\newunicodechar{⁗}{\ensuremath{''''}}
\newunicodechar{⁺}{\ensuremath{^+}}
\newunicodechar{⁻}{\ensuremath{^-}}
\newunicodechar{₋}{\ensuremath{_-}}
\newunicodechar{₌}{\ensuremath{_=}}
\newunicodechar{₍}{\ensuremath{_(}}
\newunicodechar{₎}{\ensuremath{_)}}
\newunicodechar{∂}{\ensuremath{\partial}}
\newunicodechar{∇}{\ensuremath{\nabla}}
\newunicodechar{∫}{\ensuremath{\int}}
\newunicodechar{∑}{\ensuremath{\sum}}
\newunicodechar{∏}{\ensuremath{\prod}}
\newunicodechar{√}{\ensuremath{\sqrt{}}}
\newunicodechar{∧}{\ensuremath{\land}}
\newunicodechar{∨}{\ensuremath{\lor}}
\newunicodechar{∩}{\ensuremath{\cap}}
\newunicodechar{∪}{\ensuremath{\cup}}
\newunicodechar{⊆}{\ensuremath{\subseteq}}
\newunicodechar{⊂}{\ensuremath{\subset}}
\newunicodechar{∈}{\ensuremath{\in}}
\newunicodechar{∉}{\ensuremath{\notin}}
\newunicodechar{∋}{\ensuremath{\ni}}
\newunicodechar{∅}{\ensuremath{\emptyset}}
\newunicodechar{∞}{\ensuremath{\infty}}
\newunicodechar{∃}{\ensuremath{\exists}}
\newunicodechar{α}{\ensuremath{\alpha}}
\newunicodechar{β}{\ensuremath{\beta}}
\newunicodechar{γ}{\ensuremath{\gamma}}
\newunicodechar{δ}{\ensuremath{\delta}}
\newunicodechar{ε}{\ensuremath{\epsilon}}
\newunicodechar{ζ}{\ensuremath{\zeta}}
\newunicodechar{η}{\ensuremath{\eta}}
\newunicodechar{θ}{\ensuremath{\theta}}
\newunicodechar{ι}{\ensuremath{\iota}}
\newunicodechar{κ}{\ensuremath{\kappa}}
\newunicodechar{ℓ}{\ensuremath{\ell}}
\newunicodechar{μ}{\ensuremath{\mu}}
\newunicodechar{ν}{\ensuremath{\nu}}
\newunicodechar{ξ}{\ensuremath{\xi}}
\newunicodechar{π}{\ensuremath{\pi}}
\newunicodechar{ρ}{\ensuremath{\rho}}
\newunicodechar{σ}{\ensuremath{\sigma}}
\newunicodechar{τ}{\ensuremath{\tau}}
\newunicodechar{υ}{\ensuremath{\upsilon}}
\newunicodechar{φ}{\ensuremath{\phi}}
\newunicodechar{χ}{\ensuremath{\chi}}
\newunicodechar{ψ}{\ensuremath{\psi}}
\newunicodechar{ω}{\ensuremath{\omega}}
\newunicodechar{Γ}{\ensuremath{\Gamma}}
\newunicodechar{Δ}{\ensuremath{\Delta}}
\newunicodechar{Θ}{\ensuremath{\Theta}}
\newunicodechar{Λ}{\ensuremath{\Lambda}}
\newunicodechar{Ξ}{\ensuremath{\Xi}}
\newunicodechar{Π}{\ensuremath{\Pi}}
\newunicodechar{Σ}{\ensuremath{\Sigma}}
\newunicodechar{Υ}{\ensuremath{\Upsilon}}
\newunicodechar{Φ}{\ensuremath{\Phi}}
\newunicodechar{Ψ}{\ensuremath{\Psi}}
\newunicodechar{Ω}{\ensuremath{\Omega}}

% Additional mappings for code-block characters.
\newunicodechar{⁴}{\ensuremath{{}^{4}}}
\newunicodechar{⁵}{\ensuremath{{}^{5}}}
\newunicodechar{ₑ}{\ensuremath{{}_{e}}}
\newunicodechar{ₘ}{\ensuremath{{}_{m}}}
\newunicodechar{ₚ}{\ensuremath{{}_{p}}}
\newunicodechar{ℂ}{\ensuremath{\mathbb{C}}}
\newunicodechar{↔}{\ensuremath{\leftrightarrow}}
\newunicodechar{∝}{\ensuremath{\propto}}
\newunicodechar{≠}{\ensuremath{\neq}}
\newunicodechar{≗}{\ensuremath{\circeq}}
\newunicodechar{⊔}{\ensuremath{\sqcup}}
\newunicodechar{═}{=}
\newunicodechar{□}{\ensuremath{\square}}
\newunicodechar{✓}{\ensuremath{\checkmark}}
\newunicodechar{⟩}{\ensuremath{\rangle}}
\newunicodechar{⟨}{\ensuremath{\langle}}
\newunicodechar{⟹}{\ensuremath{\Longrightarrow}}

\usepackage{setspace}
\onehalfspacing
\usepackage{amsmath,amsthm,amssymb}
\IfFileExists{agda.sty}{\usepackage{agda}}{\usepackage{latex/agda}}

% Keep Agda code inside the page as far as possible.
\renewcommand{\AgdaCodeStyle}{\fontsize{8}{9}\selectfont}
\setlength{\mathindent}{0pt}
\renewcommand{\AgdaIndentSpace}{\hspace*{0.30em}}

% Keep long literate-Agda identifiers intact without log noise.
\AtBeginEnvironment{code}{\hfuzz=400pt}
\emergencystretch=3em
\hfuzz=120pt
\hbadness=10000
\vbadness=10000

\usepackage{tikz}
\usetikzlibrary{positioning,arrows.meta,calc,backgrounds,shadows,decorations.pathreplacing,decorations.pathmorphing}

% Color palette
\definecolor{fdBlue}{RGB}{70,130,180}
\definecolor{fdRed}{RGB}{180,100,100}
\definecolor{fdGray}{RGB}{50,50,60}
\definecolor{fdLight}{RGB}{245,248,250}
\definecolor{fdAccent}{RGB}{180,100,100}

\tikzset{
    base/.style={font=\sffamily, align=center},
    stage/.style={base, rectangle, rounded corners=2mm, draw=fdBlue, thick,
                  fill=fdLight, text=fdGray, inner sep=2.6mm},
    routearrow/.style={->, >=LaTeX, thick, color=fdGray}
}

\newcommand{\AgdaFlag}[1]{\texttt{-{}-#1}}

\usepackage{hyperref}
\hypersetup{colorlinks=true,linkcolor=black,urlcolor=blue,bookmarksdepth=2}
\pdfstringdefDisableCommands{%
  \def\Sigma{Sigma}%
  \def\alpha{alpha}%
  \def\circ{o}%
  \def\cdot{.}%
  \def\leq{<=}%
  \def\mathrm#1{#1}%
  \def\mathbb#1{#1}%
  \def\texttt#1{#1}%
  \def\text#1{#1}%
  \def\ensuremath#1{#1}%
  \def\texorpdfstring#1#2{#2}%
}
\usepackage{titlesec}
\usepackage{etoolbox}
\titleformat{\chapter}[display]{\normalfont\huge\bfseries\raggedright}{\chaptertitlename\ \thechapter}{20pt}{\Huge\raggedright}
\pretocmd{\section}{\phantomsection}{}{}
\pretocmd{\subsection}{\phantomsection}{}{}
\usepackage{epigraph}
\setlength{\epigraphwidth}{0.72\textwidth}
\setcounter{secnumdepth}{0}

\begin{document}
\raggedbottom

\begin{titlepage}
\centering
\vspace*{3cm}
{\Huge\bfseries The First Distinction}\\[0.5cm]
{\Large\itshape A Companion to Void and Form}\\[2cm]
{\large Johannes Michael Wielsch}\\[4cm]
\begin{tikzpicture}[scale=2]
  \coordinate (A) at (0,0); \coordinate (B) at (1,0);
  \coordinate (C) at (0.5,0.866); \coordinate (D) at (0.5,0.289);
  \draw[thick] (A)--(B)--(C)--cycle;
  \draw[thick] (A)--(D); \draw[thick] (B)--(D); \draw[thick] (C)--(D);
  \foreach \p in {A,B,C,D} \fill (\p) circle (2pt);
\end{tikzpicture}\\[4cm]
{\small Built with AI}\\
{\small Machine-verified in Agda}\\
\vfill
{\small June 2026}\\[0.5cm]
\end{titlepage}
\clearpage

\frontmatter

\chapter*{Prologue: The Old Dream}
\addcontentsline{toc}{chapter}{Prologue: The Old Dream}

One of the oldest dreams of science is that explanation can advance by
reducing the number of independent assumptions. Leibniz imagined a reasoned
calculus strong enough to replace endless dispute with demonstrable
connection. Newton showed that the fall of a stone and the motion of the
planets need not belong to separate worlds. Maxwell brought electricity,
magnetism, and light under one field description. Einstein joined space and
time into a single kinematic structure. Modern physics has continued the
same pressure in different forms, again and again asking whether what looks
plural can be read from something more elementary.

These moments matter not only because they unify results already known. They
change what counts as basic. Domains that once appeared independent begin to
look like different expressions of a deeper order. Much of the power of
science lies precisely here: not merely in adding facts, but in discovering
that fewer principles can carry more.

This book stands inside that movement, but not at its familiar level. It
does not begin by asking which equation might join gravitation and quantum
theory, or which model might absorb all particles and forces into one more
comprehensive system. It asks a prior question.

Is there a common origin still beneath logic, mathematics, and physical
description themselves?

That question is larger than a technical unification problem and smaller
than a proclamation about reality as a whole. Nothing here claims that the
world has already been derived. The issue is whether the possibility of
description may rest on a common minimal structure.

The wager of the present project can therefore be stated plainly. Stable
information presupposes a successfully carried distinction. If that wager
fails, the route ends at once. If it holds, the question becomes how far
the first distinction can carry.

This companion takes only the first formally controlled step in that longer
route. It does not present the full theory. It does not present the
physics. It does not present the unification as accomplished. It seeks the
first point at which the conditions of stable distinction can be written
completely, constructed explicitly, and checked by machine in intensional
Martin-L\"of type theory with Agda.

Its path is therefore narrow and deliberate. Why distinction at all? What
conditions must a stable distinction satisfy? In what formal language can
those conditions be made exact? What is the first fully writable carrier of
those conditions? What already follows from so little?

Everything beyond that belongs to the larger investigation of \textit{Void}
and \textit{Form}. This book marks the threshold.

\tableofcontents

\mainmatter

% ─────────────────────────────────────────────────────────────────────
%  Literate-Agda module declaration.  The narrative runs for several
%  chapters after this minimal declaration; the formal kernel enters only
%  once the question has made it necessary.  Until then the module carries
%  only its own declaration, which keeps the source a valid literate-Agda
%  file under --safe.
% ─────────────────────────────────────────────────────────────────────

\begin{code}[hide]%
\>[0]\AgdaSymbol{\{-\#}\AgdaSpace{}%
\AgdaKeyword{OPTIONS}\AgdaSpace{}%
\AgdaPragma{--safe}\AgdaSpace{}%
\AgdaPragma{--without-K}\AgdaSpace{}%
\AgdaSymbol{\#-\}}\<%
\\
%
\\[\AgdaEmptyExtraSkip]%
\>[0]\AgdaKeyword{module}\AgdaSpace{}%
\AgdaModule{FirstDistinction}\AgdaSpace{}%
\AgdaKeyword{where}\<%
\end{code}

% ─────────────────────────────────────────────────────────────────────
\chapter{Why Distinction?}
\label{ch:origin-question}
% ─────────────────────────────────────────────────────────────────────

\epigraph{``Why is there something rather than nothing? For nothing is
simpler and easier than something.''}{Gottfried Wilhelm Leibniz}

The prologue placed this inquiry inside an older scientific movement toward
fewer independent assumptions. This chapter pushes that movement below
objects, laws, and theories themselves. Origin, determination, physics, and
information approach the matter from four different sides. They are
gathered here not as four unrelated claims, but as four witnesses to one
conjecture: that determination presupposes successful distinction.

\section{The Question of Origin}

In 1714 Leibniz formulated the question of origin in a form that remains
sharp: why is there something rather than nothing? Nothing, Leibniz says,
is simpler and easier than something. If the world could just as well have
been empty, then precisely its not being empty calls for an account.

No measurement answers this question. One can trace the history of the cosmos
back to the first physically describable state and still leave open why there
is a cosmos whose history can be traced at all. The question reaches beyond
everything data alone can decide.

But the question already carries one of its own conditions. To ask between
\emph{something} and \emph{nothing}, both must already be in force as
different. A nothing that could not be distinguished from something could
not appear as nothing. The question of origin therefore already
presupposes a distinction it does not itself explain.

That is why it leads immediately further. Not only: why is there something
rather than nothing? But also: what must already hold for something to be
distinguishable from nothing at all?

% ─────────────────────────────────────────────────────────────────────
\section{Determination}
\label{ch:description}
% ─────────────────────────────────────────────────────────────────────

A brief answer to that last question is: it must be possible at all to set
something off from something else. This condition concerns not only the
contrast between something and nothing. It holds wherever something is
determined as this and not that.

A simple case is enough.

A mark is this mark and not the blank place on the page. This sheet is this
sheet and not that one. One side is this side and not the other. Wherever
something is determinate, something else is excluded.

One cannot step behind holding-apart itself in search of a purer ground.
Without this elementary difference there would be no determinate something.
A sign that is not distinguished from what it refers to is no sign. A mark
that does not stand out from the sheet on which it stands is no mark.

At first this describes only the simplest case. But its reach is clear:
wherever something appears as this and not as that, something must already
be distinguished from something else.

If that is right, the point should show itself most clearly where
determination has been driven furthest: in physics. Does physics encounter
here only a limit of our knowledge, or a condition for anything to appear
as different at all?

% ─────────────────────────────────────────────────────────────────────
\section{The Physical Limit}
\label{ch:edge}
% ─────────────────────────────────────────────────────────────────────

If this condition is real, it must be visible in physics. There the attempt
to grasp the beginning of the universe as a state of the world has been
driven furthest.

Modern cosmology reconstructs the history of the universe backward.
Galaxies move closer together, the universe becomes hotter and denser,
atoms dissolve into plasma, nuclei break apart. This reconstruction is
quantitative, tested, and in important regions very precise. Yet it does
not reach its first state. The closer physics comes to the beginning, the
thinner the data become. Light does not reach us from before a certain
point. Energies climb beyond everything an experiment can reproduce. Space
and time themselves lose the status of magnitudes that could simply be
carried further within the same theory.

One can take this for a mere lack of data. Perhaps, so the natural
expectation goes, deeper measurements or a deeper theory will suffice. But
that expectation already assumes that better data will in the end also give
better answers about origin.

The history of physics speaks cautiously against that hope. Newton brought
celestial mechanics and falling bodies under a common law, but presupposed
space, time, and mass. Maxwell bound electricity, magnetism, and light into
a few equations, but presupposed charge and the continuum in which it
appears. Einstein made gravitation into the geometry of spacetime, but took
matter and coupling constants as given. Quantum field theory yields the most
precise predictions we have, but in turn requires a symmetry group, a
particle inventory, and constant parameters that it does not generate from
itself.

The examples share a common form. Even the most successful theories explain
more than their predecessors and yet still work with quantities,
inventories, or constants they do not generate from themselves.

The point can also be made without that history. Data are always data of
something: states, events, relations, measured values. For them to be data
at all, something must be distinguished from something else. The beginning
of the universe, however, is not simply one more state within an already
differentiated world. It concerns the question of how such a world can come
into force at all.

The first moment is therefore not only an unsolved problem of physics. It
marks the place where physics touches its own presupposition. Physics works
with distinguishable states. Remove that presupposition, and no
measurement, no dynamics, no information, and no theory remain.

The question is therefore not only: which state was the first? But: what
must hold for states to be in force at all? Then the issue is no longer
only the universe, but the condition under which anything can appear as
different at all. If states occur only as distinguished states, then
information too depends on distinction. That raises the next question:
why exactly?

% ─────────────────────────────────────────────────────────────────
\section{Information}
\label{ch:distinction-information}
% ─────────────────────────────────────────────────────────────────

Information presupposes distinction. A measured value is this value only
because it is not another. A recorded bit is this bit only because a
different outcome remains excluded.

Spencer Brown places at the beginning of his calculus the drawn
distinction. Wheeler reads physical reality as an answer to yes-or-no
questions. In both cases information hangs on retained difference.

For the rest of the inquiry, that becomes a working assumption. Every piece
of information rests on a successful act of distinction. \emph{Successful}
means: the
difference must be stateable, it must not collapse at once, and it must be
determinate enough to present itself again as the same.

The question thus shifts from a mere limit to the required performance.
What must a distinction do if it is to carry information? Can a difference
be stated at all? Does it hold? And is a field of such differences
complete, with neither less nor more?

% ─────────────────────────────────────────────────────────────────────
\chapter{Conditions of Stable Distinction}
\label{ch:distinguishability}
% ─────────────────────────────────────────────────────────────────────

The previous question asks after a performance, not a boundary. The answer
develops in three stages: a difference must be stateable, it must hold, and
an entire field of such differences must stand exactly. No one stage is
enough by itself. Only together do they yield the conditions under which a
distinction can persist stably.

\section{Distinguishability}

What is the minimum without which a distinction does not succeed?
Distinction is not a thing here but a performance: the holding-apart of one
from another and the maintenance of that difference.

A mark on a page counts as a sign only because it stands out from the page
--- because something is read \emph{as} mark and the rest \emph{as}
background. A measurement counts as a value only because it could have come
out differently: this value and not that one. In both cases the issue is
not yet counting or structure. It is the elementary possibility that
something can appear as this rather than that.

That is the first demand. Before anything can be counted, ordered, or
measured, a difference must be available. Call that distinguishability.

Nothing has yet been said about number, order, or stability.
Distinguishability is only the availability of a contrast. But it does not
suffice. A stated difference could collapse at once. What is needed so that
it does not at once fall back into one?

% ─────────────────────────────────────────────────────────────────────
\section{Non-Collapse}
\label{ch:reference}
% ─────────────────────────────────────────────────────────────────────

What distinguishability leaves open is whether a difference holds as a
difference at all. Something can lie apart for a moment and in the next
fall back into one. Then the difference carries nothing. No mark would hold
anything. No measured value could rest on it. No state could be asserted
against another.

That is why a second demand follows distinguishability. A difference does
not suffice merely because it flashes up once. It must be secured against
collapse. Call this demand Non-Collapse.

Long before there are scales, observers, or instruments, it must be
possible for a difference not to vanish at once. Measurement does not
create that. Recording does not create it. Both already presuppose it.

The same demand appears in physics as well, for example in the black-hole
information problem. The issue there is not first of matter or energy, but
of whether differences once achieved can be lost without remainder. This is
not a proof of the present argument, but it is a clear physical case of the
same structure.

The demand can now be sharpened. Non-Collapse means: there must be a genuine
separation between two positions, and that separation may not fall back into
identity. That is exactly why the formal part will later have to work with
inequality and equality.

With this much in hand, the next question arises. Is a single difference
secured against collapse enough, or must an entire field of such
differences stand complete?

% ─────────────────────────────────────────────────────────────────────
\section{Exact Plurality}
\label{ch:plurality}
% ─────────────────────────────────────────────────────────────────────

The next question no longer concerns a single difference, but an entire
field of differences. What does it take for a system to be fixed to a
determinate plurality: that exactly these positions are here, neither fewer
nor more?

This demand is felt most directly by trying to hold a small plurality in
place and seeing what can go wrong. Consider a few blank sheets on a table
and suppose there is a definite number of them. Two failures lie in wait,
and they pull in opposite directions. The first: some of the sheets counted
as different might not really be different at all --- if every way of
pointing to one points equally well to another, then there are not distinct
sheets, but one sheet pointed at twice. Call this collapse: the many
secretly fall together into fewer. The second, less obvious case: there may
be more on the table than the count ever mentioned --- a sheet at the edge
that the accounting never ruled out. Call this surplus: the field hides
elements the count never closed.

These two dangers are distinct. Collapse means fewer than claimed. Surplus
means more than claimed. Whoever prevents only collapse may still have too
much in the field. Whoever excludes only surplus may still secretly identify
two positions. An exact plurality has to answer both dangers at once.

The demand that answers both has exactly two parts. The named positions must
really lie apart. Call that \emph{separation}. And every position of the
field must be captured by the account. Call that \emph{cover}. Separation
prevents collapse. Cover prevents surplus. Only both together turn named
positions into an exact plurality.

With separation and cover we have, for the first time, a specification
complete enough to be written down exactly.

Ask where these two conditions can first be met completely, and the answer
is: on the smallest field on which anything can be held apart at all,
namely a field of two positions. With two, and not before, separation and
cover can be written down completely and finitely. Two is not the origin of
distinctness. Two is the first place at which its general conditions become
fully explicit.

If this is the minimal performance stable information requires, how little
must be granted to write it down exactly?

% ─────────────────────────────────────────────────────────────────
\chapter{The Threshold of Formalization}
\label{ch:crossing}
% ─────────────────────────────────────────────────────────────────

At this point the question changes character.

Up to now the discussion has been exploratory. The issue was whether certain
conditions seem necessary for a distinction to persist. But necessity stated
in ordinary language is still only a proposal.

The next question is stricter.

Can these conditions be written in a form precise enough that every
consequence becomes checkable and every hidden assumption becomes visible?

The inquiry therefore moves from philosophical orientation to formal
specification.

From here on, every claim must be paid for by construction.

The previous chapter narrowed the question to a small finite
specification: separation and cover on a field of two positions. At that
point ordinary language reaches its limit. What is now missing is not
another argument, but a language in which every step can be stated exactly
and checked.

There is such a language. The one used here is Martin-L\"of type theory in
the form realized by the proof assistant Agda.

What is granted can be said briefly: a carrier with two positions, a
witness that the two are not equal, and an account that leaves no point of
the carrier unaddressed. No numbers, no geometry, no physics, and no
classical logic beyond what is constructed explicitly.

Everything else must be built.

In this setting, to assert is to construct. A statement is a type, a proof
is a term inhabiting that type, and a machine decides whether the term
truly has the type claimed. A proof that does not check is simply a file
that does not compile.

The argument is formulated in type theory not because it is the only
possible carrier, but because it makes the cost of every step visible.

Two settings, fixed at the head of the file, govern everything that
follows.

The flag \AgdaFlag{safe} closes every escape hatch: no postulates, no
principle asserted without proof, no back door through which an unearned
result could enter.

The flag \AgdaFlag{without-K} withholds a convenient but additional
assumption about equality, so that nothing in the development rests on more
than the bare identity type.

Under these restrictions, only the things we actually build exist.

If the formalization later seems to do more than this brief list suggests,
that is only because much already follows from little. It is not because
more was smuggled in beforehand. The order of writing is not the same as 
the order of the inquiry.

In the inquiry, distinction, Non-Collapse, and exact plurality appeared
first. In the file, what appears first is whatever the checker requires in
order for those conditions to be sayable at all.

No new content is introduced here. The file merely provides the few means
without which separation and cover could not be stated exactly. These 
constructions belong to the requirements of the formal language, not
to the order in which the inquiry itself arrived at its result.

It therefore begins with the bookkeeping required to manage universe
levels. Nothing conceptual enters with that step.

\begin{code}%
\>[0]\AgdaKeyword{open}\AgdaSpace{}%
\AgdaKeyword{import}\AgdaSpace{}%
\AgdaModule{Agda.Primitive}\AgdaSpace{}%
\AgdaKeyword{using}\AgdaSpace{}%
\AgdaSymbol{(}\AgdaPostulate{Level}\AgdaSymbol{;}\AgdaSpace{}%
\AgdaPrimitive{lzero}\AgdaSymbol{;}\AgdaSpace{}%
\AgdaPrimitive{lsuc}\AgdaSymbol{;}\AgdaSpace{}%
\AgdaOperator{\AgdaPrimitive{\AgdaUnderscore{}⊔\AgdaUnderscore{}}}\AgdaSymbol{;}\AgdaSpace{}%
\AgdaPrimitive{Setω}\AgdaSymbol{)}\<%
\end{code}

That first line is only housekeeping. It gives the checker a way to track
sizes consistently, and none of the conceptual weight of the argument rests
on it.

The real formal work begins with the thinnest device needed to say when
something is the same and when it is not.

That is why equality comes first in the file. Not because sameness would be
prior to distinction in the order of the inquiry, but because separation
and cover are not even formulable without a minimal way to say same, not
the same, and substitution.

The identity type provides exactly that, and at this stage it is not meant
to do more.

\begin{code}%
\>[0]\AgdaKeyword{infix}\AgdaSpace{}%
\AgdaNumber{4}\AgdaSpace{}%
\AgdaOperator{\AgdaDatatype{\AgdaUnderscore{}≡\AgdaUnderscore{}}}\<%
\\
%
\\[\AgdaEmptyExtraSkip]%
\>[0]\AgdaKeyword{data}\AgdaSpace{}%
\AgdaOperator{\AgdaDatatype{\AgdaUnderscore{}≡\AgdaUnderscore{}}}\AgdaSpace{}%
\AgdaSymbol{\{}\AgdaBound{ℓ}\AgdaSpace{}%
\AgdaSymbol{:}\AgdaSpace{}%
\AgdaPostulate{Level}\AgdaSymbol{\}}\AgdaSpace{}%
\AgdaSymbol{\{}\AgdaBound{A}\AgdaSpace{}%
\AgdaSymbol{:}\AgdaSpace{}%
\AgdaPrimitive{Set}\AgdaSpace{}%
\AgdaBound{ℓ}\AgdaSymbol{\}}\AgdaSpace{}%
\AgdaSymbol{(}\AgdaBound{x}\AgdaSpace{}%
\AgdaSymbol{:}\AgdaSpace{}%
\AgdaBound{A}\AgdaSymbol{)}\AgdaSpace{}%
\AgdaSymbol{:}\AgdaSpace{}%
\AgdaBound{A}\AgdaSpace{}%
\AgdaSymbol{→}\AgdaSpace{}%
\AgdaPrimitive{Set}\AgdaSpace{}%
\AgdaBound{ℓ}\AgdaSpace{}%
\AgdaKeyword{where}\<%
\\
\>[0][@{}l@{\AgdaIndent{0}}]%
\>[2]\AgdaInductiveConstructor{refl}\AgdaSpace{}%
\AgdaSymbol{:}\AgdaSpace{}%
\AgdaBound{x}\AgdaSpace{}%
\AgdaOperator{\AgdaDatatype{≡}}\AgdaSpace{}%
\AgdaBound{x}\<%
\\
\>[0]\AgdaSymbol{\{-\#}\AgdaSpace{}%
\AgdaKeyword{BUILTIN}\AgdaSpace{}%
\AgdaKeyword{EQUALITY}\AgdaSpace{}%
\AgdaOperator{\AgdaDatatype{\AgdaUnderscore{}≡\AgdaUnderscore{}}}\AgdaSpace{}%
\AgdaSymbol{\#-\}}\<%
\\
%
\\[\AgdaEmptyExtraSkip]%
\>[0]\AgdaFunction{sym}\AgdaSpace{}%
\AgdaSymbol{:}\AgdaSpace{}%
\AgdaSymbol{\{}\AgdaBound{ℓ}\AgdaSpace{}%
\AgdaSymbol{:}\AgdaSpace{}%
\AgdaPostulate{Level}\AgdaSymbol{\}}\AgdaSpace{}%
\AgdaSymbol{\{}\AgdaBound{A}\AgdaSpace{}%
\AgdaSymbol{:}\AgdaSpace{}%
\AgdaPrimitive{Set}\AgdaSpace{}%
\AgdaBound{ℓ}\AgdaSymbol{\}}\AgdaSpace{}%
\AgdaSymbol{\{}\AgdaBound{x}\AgdaSpace{}%
\AgdaBound{y}\AgdaSpace{}%
\AgdaSymbol{:}\AgdaSpace{}%
\AgdaBound{A}\AgdaSymbol{\}}\AgdaSpace{}%
\AgdaSymbol{→}\AgdaSpace{}%
\AgdaBound{x}\AgdaSpace{}%
\AgdaOperator{\AgdaDatatype{≡}}\AgdaSpace{}%
\AgdaBound{y}\AgdaSpace{}%
\AgdaSymbol{→}\AgdaSpace{}%
\AgdaBound{y}\AgdaSpace{}%
\AgdaOperator{\AgdaDatatype{≡}}\AgdaSpace{}%
\AgdaBound{x}\<%
\\
\>[0]\AgdaFunction{sym}\AgdaSpace{}%
\AgdaInductiveConstructor{refl}\AgdaSpace{}%
\AgdaSymbol{=}\AgdaSpace{}%
\AgdaInductiveConstructor{refl}\<%
\\
%
\\[\AgdaEmptyExtraSkip]%
\>[0]\AgdaFunction{trans}\AgdaSpace{}%
\AgdaSymbol{:}\AgdaSpace{}%
\AgdaSymbol{\{}\AgdaBound{ℓ}\AgdaSpace{}%
\AgdaSymbol{:}\AgdaSpace{}%
\AgdaPostulate{Level}\AgdaSymbol{\}}\AgdaSpace{}%
\AgdaSymbol{\{}\AgdaBound{A}\AgdaSpace{}%
\AgdaSymbol{:}\AgdaSpace{}%
\AgdaPrimitive{Set}\AgdaSpace{}%
\AgdaBound{ℓ}\AgdaSymbol{\}}\AgdaSpace{}%
\AgdaSymbol{\{}\AgdaBound{x}\AgdaSpace{}%
\AgdaBound{y}\AgdaSpace{}%
\AgdaBound{z}\AgdaSpace{}%
\AgdaSymbol{:}\AgdaSpace{}%
\AgdaBound{A}\AgdaSymbol{\}}\AgdaSpace{}%
\AgdaSymbol{→}\AgdaSpace{}%
\AgdaBound{x}\AgdaSpace{}%
\AgdaOperator{\AgdaDatatype{≡}}\AgdaSpace{}%
\AgdaBound{y}\AgdaSpace{}%
\AgdaSymbol{→}\AgdaSpace{}%
\AgdaBound{y}\AgdaSpace{}%
\AgdaOperator{\AgdaDatatype{≡}}\AgdaSpace{}%
\AgdaBound{z}\AgdaSpace{}%
\AgdaSymbol{→}\AgdaSpace{}%
\AgdaBound{x}\AgdaSpace{}%
\AgdaOperator{\AgdaDatatype{≡}}\AgdaSpace{}%
\AgdaBound{z}\<%
\\
\>[0]\AgdaFunction{trans}\AgdaSpace{}%
\AgdaInductiveConstructor{refl}\AgdaSpace{}%
\AgdaBound{q}\AgdaSpace{}%
\AgdaSymbol{=}\AgdaSpace{}%
\AgdaBound{q}\<%
\\
%
\\[\AgdaEmptyExtraSkip]%
\>[0]\AgdaFunction{cong}\AgdaSpace{}%
\AgdaSymbol{:}\AgdaSpace{}%
\AgdaSymbol{\{}\AgdaBound{ℓ₁}\AgdaSpace{}%
\AgdaBound{ℓ₂}\AgdaSpace{}%
\AgdaSymbol{:}\AgdaSpace{}%
\AgdaPostulate{Level}\AgdaSymbol{\}}\AgdaSpace{}%
\AgdaSymbol{\{}\AgdaBound{A}\AgdaSpace{}%
\AgdaSymbol{:}\AgdaSpace{}%
\AgdaPrimitive{Set}\AgdaSpace{}%
\AgdaBound{ℓ₁}\AgdaSymbol{\}}\AgdaSpace{}%
\AgdaSymbol{\{}\AgdaBound{B}\AgdaSpace{}%
\AgdaSymbol{:}\AgdaSpace{}%
\AgdaPrimitive{Set}\AgdaSpace{}%
\AgdaBound{ℓ₂}\AgdaSymbol{\}}\AgdaSpace{}%
\AgdaSymbol{(}\AgdaBound{f}\AgdaSpace{}%
\AgdaSymbol{:}\AgdaSpace{}%
\AgdaBound{A}\AgdaSpace{}%
\AgdaSymbol{→}\AgdaSpace{}%
\AgdaBound{B}\AgdaSymbol{)}\AgdaSpace{}%
\AgdaSymbol{\{}\AgdaBound{x}\AgdaSpace{}%
\AgdaBound{y}\AgdaSpace{}%
\AgdaSymbol{:}\AgdaSpace{}%
\AgdaBound{A}\AgdaSymbol{\}}\AgdaSpace{}%
\AgdaSymbol{→}\AgdaSpace{}%
\AgdaBound{x}\AgdaSpace{}%
\AgdaOperator{\AgdaDatatype{≡}}\AgdaSpace{}%
\AgdaBound{y}\AgdaSpace{}%
\AgdaSymbol{→}\AgdaSpace{}%
\AgdaBound{f}\AgdaSpace{}%
\AgdaBound{x}\AgdaSpace{}%
\AgdaOperator{\AgdaDatatype{≡}}\AgdaSpace{}%
\AgdaBound{f}\AgdaSpace{}%
\AgdaBound{y}\<%
\\
\>[0]\AgdaFunction{cong}\AgdaSpace{}%
\AgdaBound{f}\AgdaSpace{}%
\AgdaInductiveConstructor{refl}\AgdaSpace{}%
\AgdaSymbol{=}\AgdaSpace{}%
\AgdaInductiveConstructor{refl}\<%
\\
%
\\[\AgdaEmptyExtraSkip]%
\>[0]\AgdaFunction{subst}\AgdaSpace{}%
\AgdaSymbol{:}\AgdaSpace{}%
\AgdaSymbol{\{}\AgdaBound{ℓ₁}\AgdaSpace{}%
\AgdaBound{ℓ₂}\AgdaSpace{}%
\AgdaSymbol{:}\AgdaSpace{}%
\AgdaPostulate{Level}\AgdaSymbol{\}}\AgdaSpace{}%
\AgdaSymbol{\{}\AgdaBound{A}\AgdaSpace{}%
\AgdaSymbol{:}\AgdaSpace{}%
\AgdaPrimitive{Set}\AgdaSpace{}%
\AgdaBound{ℓ₁}\AgdaSymbol{\}}\AgdaSpace{}%
\AgdaSymbol{(}\AgdaBound{P}\AgdaSpace{}%
\AgdaSymbol{:}\AgdaSpace{}%
\AgdaBound{A}\AgdaSpace{}%
\AgdaSymbol{→}\AgdaSpace{}%
\AgdaPrimitive{Set}\AgdaSpace{}%
\AgdaBound{ℓ₂}\AgdaSymbol{)}\AgdaSpace{}%
\AgdaSymbol{\{}\AgdaBound{x}\AgdaSpace{}%
\AgdaBound{y}\AgdaSpace{}%
\AgdaSymbol{:}\AgdaSpace{}%
\AgdaBound{A}\AgdaSymbol{\}}\AgdaSpace{}%
\AgdaSymbol{→}\AgdaSpace{}%
\AgdaBound{x}\AgdaSpace{}%
\AgdaOperator{\AgdaDatatype{≡}}\AgdaSpace{}%
\AgdaBound{y}\AgdaSpace{}%
\AgdaSymbol{→}\AgdaSpace{}%
\AgdaBound{P}\AgdaSpace{}%
\AgdaBound{x}\AgdaSpace{}%
\AgdaSymbol{→}\AgdaSpace{}%
\AgdaBound{P}\AgdaSpace{}%
\AgdaBound{y}\<%
\\
\>[0]\AgdaFunction{subst}\AgdaSpace{}%
\AgdaBound{P}\AgdaSpace{}%
\AgdaInductiveConstructor{refl}\AgdaSpace{}%
\AgdaBound{px}\AgdaSpace{}%
\AgdaSymbol{=}\AgdaSpace{}%
\AgdaBound{px}\<%
\end{code}

Next the file needs a way to say impossibility. A guard against collapse is
worthless unless one can express that an identification is excluded. The
empty type gives that means. Negation is then the map into impossibility,
and inequality --- the formal shape of being kept apart --- is the
negation of equality.

\begin{code}%
\>[0]\AgdaKeyword{data}\AgdaSpace{}%
\AgdaDatatype{⊥}\AgdaSpace{}%
\AgdaSymbol{:}\AgdaSpace{}%
\AgdaPrimitive{Set}\AgdaSpace{}%
\AgdaKeyword{where}\<%
\\
%
\\[\AgdaEmptyExtraSkip]%
\>[0]\AgdaFunction{⊥-elim}\AgdaSpace{}%
\AgdaSymbol{:}\AgdaSpace{}%
\AgdaSymbol{\{}\AgdaBound{ℓ}\AgdaSpace{}%
\AgdaSymbol{:}\AgdaSpace{}%
\AgdaPostulate{Level}\AgdaSymbol{\}}\AgdaSpace{}%
\AgdaSymbol{\{}\AgdaBound{A}\AgdaSpace{}%
\AgdaSymbol{:}\AgdaSpace{}%
\AgdaPrimitive{Set}\AgdaSpace{}%
\AgdaBound{ℓ}\AgdaSymbol{\}}\AgdaSpace{}%
\AgdaSymbol{→}\AgdaSpace{}%
\AgdaDatatype{⊥}\AgdaSpace{}%
\AgdaSymbol{→}\AgdaSpace{}%
\AgdaBound{A}\<%
\\
\>[0]\AgdaFunction{⊥-elim}\AgdaSpace{}%
\AgdaSymbol{()}\<%
\\
%
\\[\AgdaEmptyExtraSkip]%
\>[0]\AgdaOperator{\AgdaFunction{¬\AgdaUnderscore{}}}\AgdaSpace{}%
\AgdaSymbol{:}\AgdaSpace{}%
\AgdaSymbol{\{}\AgdaBound{ℓ}\AgdaSpace{}%
\AgdaSymbol{:}\AgdaSpace{}%
\AgdaPostulate{Level}\AgdaSymbol{\}}\AgdaSpace{}%
\AgdaSymbol{→}\AgdaSpace{}%
\AgdaPrimitive{Set}\AgdaSpace{}%
\AgdaBound{ℓ}\AgdaSpace{}%
\AgdaSymbol{→}\AgdaSpace{}%
\AgdaPrimitive{Set}\AgdaSpace{}%
\AgdaBound{ℓ}\<%
\\
\>[0]\AgdaOperator{\AgdaFunction{¬}}\AgdaSpace{}%
\AgdaBound{A}\AgdaSpace{}%
\AgdaSymbol{=}\AgdaSpace{}%
\AgdaBound{A}\AgdaSpace{}%
\AgdaSymbol{→}\AgdaSpace{}%
\AgdaDatatype{⊥}\<%
\\
%
\\[\AgdaEmptyExtraSkip]%
\>[0]\AgdaKeyword{infix}\AgdaSpace{}%
\AgdaNumber{2}\AgdaSpace{}%
\AgdaOperator{\AgdaFunction{\AgdaUnderscore{}≠\AgdaUnderscore{}}}\<%
\\
\>[0]\AgdaOperator{\AgdaFunction{\AgdaUnderscore{}≠\AgdaUnderscore{}}}\AgdaSpace{}%
\AgdaSymbol{:}\AgdaSpace{}%
\AgdaSymbol{\{}\AgdaBound{ℓ}\AgdaSpace{}%
\AgdaSymbol{:}\AgdaSpace{}%
\AgdaPostulate{Level}\AgdaSymbol{\}}\AgdaSpace{}%
\AgdaSymbol{\{}\AgdaBound{A}\AgdaSpace{}%
\AgdaSymbol{:}\AgdaSpace{}%
\AgdaPrimitive{Set}\AgdaSpace{}%
\AgdaBound{ℓ}\AgdaSymbol{\}}\AgdaSpace{}%
\AgdaSymbol{→}\AgdaSpace{}%
\AgdaBound{A}\AgdaSpace{}%
\AgdaSymbol{→}\AgdaSpace{}%
\AgdaBound{A}\AgdaSpace{}%
\AgdaSymbol{→}\AgdaSpace{}%
\AgdaPrimitive{Set}\AgdaSpace{}%
\AgdaBound{ℓ}\<%
\\
\>[0]\AgdaBound{x}\AgdaSpace{}%
\AgdaOperator{\AgdaFunction{≠}}\AgdaSpace{}%
\AgdaBound{y}\AgdaSpace{}%
\AgdaSymbol{=}\AgdaSpace{}%
\AgdaOperator{\AgdaFunction{¬}}\AgdaSpace{}%
\AgdaSymbol{(}\AgdaBound{x}\AgdaSpace{}%
\AgdaOperator{\AgdaDatatype{≡}}\AgdaSpace{}%
\AgdaBound{y}\AgdaSymbol{)}\<%
\end{code}

Then the file needs a way to say exactly: one of two, and which one. Cover
cannot be expressed by an untagged choice, because that would merge the two
cases. The disjoint union keeps the mark, so that a later case analysis can
say from which side a witness comes.

\begin{code}%
\>[0]\AgdaKeyword{infixr}\AgdaSpace{}%
\AgdaNumber{2}\AgdaSpace{}%
\AgdaOperator{\AgdaDatatype{\AgdaUnderscore{}⊎\AgdaUnderscore{}}}\<%
\\
%
\\[\AgdaEmptyExtraSkip]%
\>[0]\AgdaKeyword{data}\AgdaSpace{}%
\AgdaOperator{\AgdaDatatype{\AgdaUnderscore{}⊎\AgdaUnderscore{}}}\AgdaSpace{}%
\AgdaSymbol{\{}\AgdaBound{ℓ}\AgdaSpace{}%
\AgdaSymbol{:}\AgdaSpace{}%
\AgdaPostulate{Level}\AgdaSymbol{\}}\AgdaSpace{}%
\AgdaSymbol{(}\AgdaBound{A}\AgdaSpace{}%
\AgdaBound{B}\AgdaSpace{}%
\AgdaSymbol{:}\AgdaSpace{}%
\AgdaPrimitive{Set}\AgdaSpace{}%
\AgdaBound{ℓ}\AgdaSymbol{)}\AgdaSpace{}%
\AgdaSymbol{:}\AgdaSpace{}%
\AgdaPrimitive{Set}\AgdaSpace{}%
\AgdaBound{ℓ}\AgdaSpace{}%
\AgdaKeyword{where}\<%
\\
\>[0][@{}l@{\AgdaIndent{0}}]%
\>[2]\AgdaInductiveConstructor{inj₁}\AgdaSpace{}%
\AgdaSymbol{:}\AgdaSpace{}%
\AgdaBound{A}\AgdaSpace{}%
\AgdaSymbol{→}\AgdaSpace{}%
\AgdaBound{A}\AgdaSpace{}%
\AgdaOperator{\AgdaDatatype{⊎}}\AgdaSpace{}%
\AgdaBound{B}\<%
\\
%
\>[2]\AgdaInductiveConstructor{inj₂}\AgdaSpace{}%
\AgdaSymbol{:}\AgdaSpace{}%
\AgdaBound{B}\AgdaSpace{}%
\AgdaSymbol{→}\AgdaSpace{}%
\AgdaBound{A}\AgdaSpace{}%
\AgdaOperator{\AgdaDatatype{⊎}}\AgdaSpace{}%
\AgdaBound{B}\<%
\end{code}

Finally, the file needs a way to carry a thing together with a statement
about it. That is what the dependent pair does, with the ordinary product
as its constant special case. In this way an existential claim becomes an
actual witness together with a proof about it.

\begin{code}%
\>[0]\AgdaKeyword{infixr}\AgdaSpace{}%
\AgdaNumber{3}\AgdaSpace{}%
\AgdaOperator{\AgdaRecord{\AgdaUnderscore{}×\AgdaUnderscore{}}}\<%
\\
%
\\[\AgdaEmptyExtraSkip]%
\>[0]\AgdaKeyword{record}\AgdaSpace{}%
\AgdaOperator{\AgdaRecord{\AgdaUnderscore{}×\AgdaUnderscore{}}}\AgdaSpace{}%
\AgdaSymbol{(}\AgdaBound{A}\AgdaSpace{}%
\AgdaBound{B}\AgdaSpace{}%
\AgdaSymbol{:}\AgdaSpace{}%
\AgdaPrimitive{Set}\AgdaSymbol{)}\AgdaSpace{}%
\AgdaSymbol{:}\AgdaSpace{}%
\AgdaPrimitive{Set}\AgdaSpace{}%
\AgdaKeyword{where}\<%
\\
\>[0][@{}l@{\AgdaIndent{0}}]%
\>[2]\AgdaKeyword{constructor}\AgdaSpace{}%
\AgdaOperator{\AgdaInductiveConstructor{\AgdaUnderscore{},\AgdaUnderscore{}}}\<%
\\
%
\>[2]\AgdaKeyword{field}\<%
\\
\>[2][@{}l@{\AgdaIndent{0}}]%
\>[4]\AgdaField{fst}\AgdaSpace{}%
\AgdaSymbol{:}\AgdaSpace{}%
\AgdaBound{A}\<%
\\
%
\>[4]\AgdaField{snd}\AgdaSpace{}%
\AgdaSymbol{:}\AgdaSpace{}%
\AgdaBound{B}\<%
\\
%
\\[\AgdaEmptyExtraSkip]%
\>[0]\AgdaKeyword{open}\AgdaSpace{}%
\AgdaOperator{\AgdaModule{\AgdaUnderscore{}×\AgdaUnderscore{}}}\AgdaSpace{}%
\AgdaKeyword{public}\<%
\\
%
\\[\AgdaEmptyExtraSkip]%
\>[0]\AgdaKeyword{record}\AgdaSpace{}%
\AgdaRecord{Σ}\AgdaSpace{}%
\AgdaSymbol{(}\AgdaBound{A}\AgdaSpace{}%
\AgdaSymbol{:}\AgdaSpace{}%
\AgdaPrimitive{Set}\AgdaSymbol{)}\AgdaSpace{}%
\AgdaSymbol{(}\AgdaBound{B}\AgdaSpace{}%
\AgdaSymbol{:}\AgdaSpace{}%
\AgdaBound{A}\AgdaSpace{}%
\AgdaSymbol{→}\AgdaSpace{}%
\AgdaPrimitive{Set}\AgdaSymbol{)}\AgdaSpace{}%
\AgdaSymbol{:}\AgdaSpace{}%
\AgdaPrimitive{Set}\AgdaSpace{}%
\AgdaKeyword{where}\<%
\\
\>[0][@{}l@{\AgdaIndent{0}}]%
\>[2]\AgdaKeyword{constructor}\AgdaSpace{}%
\AgdaOperator{\AgdaInductiveConstructor{\AgdaUnderscore{},\AgdaUnderscore{}}}\<%
\\
%
\>[2]\AgdaKeyword{field}\<%
\\
\>[2][@{}l@{\AgdaIndent{0}}]%
\>[4]\AgdaField{fst}\AgdaSpace{}%
\AgdaSymbol{:}\AgdaSpace{}%
\AgdaBound{A}\<%
\\
%
\>[4]\AgdaField{snd}\AgdaSpace{}%
\AgdaSymbol{:}\AgdaSpace{}%
\AgdaBound{B}\AgdaSpace{}%
\AgdaField{fst}\<%
\\
%
\\[\AgdaEmptyExtraSkip]%
\>[0]\AgdaKeyword{open}\AgdaSpace{}%
\AgdaModule{Σ}\AgdaSpace{}%
\AgdaKeyword{public}\<%
\end{code}

With these means in place, the earlier question can for the first time be
translated into a single exact object. What does this smallest object look
like, the one that carries separation and cover together?

% ─────────────────────────────────────────────────────────────────
\chapter{The First Distinction}
\label{ch:the-record}
% ─────────────────────────────────────────────────────────────────

The previous question asked how little has to be granted for the minimal
conditions of successful distinction to be writable exactly. Here the title
of the book is earned. The Distinction Record is the first complete answer.
Formally it is a record. In content it is the first fully writable carrier
of a stable distinction: not an inventory of a thing, but a certificate of
conditions, exactly the data needed for a difference to persist stably as a
difference.

\begin{code}%
\>[0]\AgdaKeyword{record}\AgdaSpace{}%
\AgdaRecord{Distinction}\AgdaSpace{}%
\AgdaSymbol{:}\AgdaSpace{}%
\AgdaPrimitive{Set1}\AgdaSpace{}%
\AgdaKeyword{where}\<%
\\
\>[0][@{}l@{\AgdaIndent{0}}]%
\>[2]\AgdaKeyword{field}\<%
\\
\>[2][@{}l@{\AgdaIndent{0}}]%
\>[4]\AgdaField{S}%
\>[10]\AgdaSymbol{:}\AgdaSpace{}%
\AgdaPrimitive{Set}\<%
\\
%
\>[4]\AgdaField{ℓ}%
\>[10]\AgdaSymbol{:}\AgdaSpace{}%
\AgdaField{S}\<%
\\
%
\>[4]\AgdaField{r}%
\>[10]\AgdaSymbol{:}\AgdaSpace{}%
\AgdaField{S}\<%
\\
%
\>[4]\AgdaField{ℓ≠r}%
\>[10]\AgdaSymbol{:}\AgdaSpace{}%
\AgdaField{ℓ}\AgdaSpace{}%
\AgdaOperator{\AgdaFunction{≠}}\AgdaSpace{}%
\AgdaField{r}\<%
\\
%
\>[4]\AgdaField{cover}\AgdaSpace{}%
\AgdaSymbol{:}\AgdaSpace{}%
\AgdaSymbol{(}\AgdaBound{x}\AgdaSpace{}%
\AgdaSymbol{:}\AgdaSpace{}%
\AgdaField{S}\AgdaSymbol{)}\AgdaSpace{}%
\AgdaSymbol{→}\AgdaSpace{}%
\AgdaSymbol{(}\AgdaBound{x}\AgdaSpace{}%
\AgdaOperator{\AgdaDatatype{≡}}\AgdaSpace{}%
\AgdaField{ℓ}\AgdaSymbol{)}\AgdaSpace{}%
\AgdaOperator{\AgdaDatatype{⊎}}\AgdaSpace{}%
\AgdaSymbol{(}\AgdaBound{x}\AgdaSpace{}%
\AgdaOperator{\AgdaDatatype{≡}}\AgdaSpace{}%
\AgdaField{r}\AgdaSymbol{)}\<%
\\
%
\\[\AgdaEmptyExtraSkip]%
\>[0]\AgdaKeyword{open}\AgdaSpace{}%
\AgdaModule{Distinction}\AgdaSpace{}%
\AgdaKeyword{public}\<%
\end{code}

Read back into prose, \texttt{S} is the carrier, \texttt{ℓ} and
\texttt{r} are the two positions, \texttt{ℓ≠r} is separation, and
\texttt{cover} is cover. The record contains no more than that. That is
its point. It does not primarily describe an object. It describes the
minimal configuration under which a distinction can be held fast as a
successful distinction. Its first consequence is simple: \texttt{cover}
says that every point of the carrier reduces to \texttt{ℓ} or \texttt{r}.

\begin{code}%
\>[0]\AgdaFunction{law1-1-cover}\AgdaSpace{}%
\AgdaSymbol{:}\AgdaSpace{}%
\AgdaSymbol{(}\AgdaBound{d}\AgdaSpace{}%
\AgdaSymbol{:}\AgdaSpace{}%
\AgdaRecord{Distinction}\AgdaSymbol{)}\AgdaSpace{}%
\AgdaSymbol{→}\AgdaSpace{}%
\AgdaSymbol{(}\AgdaBound{x}\AgdaSpace{}%
\AgdaSymbol{:}\AgdaSpace{}%
\AgdaField{S}\AgdaSpace{}%
\AgdaBound{d}\AgdaSymbol{)}\AgdaSpace{}%
\AgdaSymbol{→}\AgdaSpace{}%
\AgdaSymbol{(}\AgdaBound{x}\AgdaSpace{}%
\AgdaOperator{\AgdaDatatype{≡}}\AgdaSpace{}%
\AgdaField{ℓ}\AgdaSpace{}%
\AgdaBound{d}\AgdaSymbol{)}\AgdaSpace{}%
\AgdaOperator{\AgdaDatatype{⊎}}\AgdaSpace{}%
\AgdaSymbol{(}\AgdaBound{x}\AgdaSpace{}%
\AgdaOperator{\AgdaDatatype{≡}}\AgdaSpace{}%
\AgdaField{r}\AgdaSpace{}%
\AgdaBound{d}\AgdaSymbol{)}\<%
\\
\>[0]\AgdaFunction{law1-1-cover}\AgdaSpace{}%
\AgdaSymbol{=}\AgdaSpace{}%
\AgdaField{cover}\<%
\end{code}

If every point is one of the two, then it is enough to check a property at
\texttt{ℓ} and \texttt{r} in order to obtain it for all points. That is
the carrier eliminator.

\begin{code}%
\>[0]\AgdaFunction{Distinction-elim}\AgdaSpace{}%
\AgdaSymbol{:}\<%
\\
\>[0][@{}l@{\AgdaIndent{0}}]%
\>[2]\AgdaSymbol{(}\AgdaBound{d}\AgdaSpace{}%
\AgdaSymbol{:}\AgdaSpace{}%
\AgdaRecord{Distinction}\AgdaSymbol{)}\AgdaSpace{}%
\AgdaSymbol{→}\<%
\\
%
\>[2]\AgdaSymbol{\{}\AgdaBound{P}\AgdaSpace{}%
\AgdaSymbol{:}\AgdaSpace{}%
\AgdaField{S}\AgdaSpace{}%
\AgdaBound{d}\AgdaSpace{}%
\AgdaSymbol{→}\AgdaSpace{}%
\AgdaPrimitive{Set}\AgdaSymbol{\}}\AgdaSpace{}%
\AgdaSymbol{→}\<%
\\
%
\>[2]\AgdaBound{P}\AgdaSpace{}%
\AgdaSymbol{(}\AgdaField{ℓ}\AgdaSpace{}%
\AgdaBound{d}\AgdaSymbol{)}\AgdaSpace{}%
\AgdaSymbol{→}\<%
\\
%
\>[2]\AgdaBound{P}\AgdaSpace{}%
\AgdaSymbol{(}\AgdaField{r}\AgdaSpace{}%
\AgdaBound{d}\AgdaSymbol{)}\AgdaSpace{}%
\AgdaSymbol{→}\<%
\\
%
\>[2]\AgdaSymbol{(}\AgdaBound{x}\AgdaSpace{}%
\AgdaSymbol{:}\AgdaSpace{}%
\AgdaField{S}\AgdaSpace{}%
\AgdaBound{d}\AgdaSymbol{)}\AgdaSpace{}%
\AgdaSymbol{→}\<%
\\
%
\>[2]\AgdaBound{P}\AgdaSpace{}%
\AgdaBound{x}\<%
\\
\>[0]\AgdaFunction{Distinction-elim}\AgdaSpace{}%
\AgdaBound{d}\AgdaSpace{}%
\AgdaSymbol{\{}\AgdaBound{P}\AgdaSymbol{\}}\AgdaSpace{}%
\AgdaBound{pℓ}\AgdaSpace{}%
\AgdaBound{pr}\AgdaSpace{}%
\AgdaBound{x}\AgdaSpace{}%
\AgdaKeyword{with}\AgdaSpace{}%
\AgdaField{cover}\AgdaSpace{}%
\AgdaBound{d}\AgdaSpace{}%
\AgdaBound{x}\<%
\\
\>[0]\AgdaSymbol{...}\AgdaSpace{}%
\AgdaSymbol{|}\AgdaSpace{}%
\AgdaInductiveConstructor{inj₁}\AgdaSpace{}%
\AgdaBound{x≡ℓ}\AgdaSpace{}%
\AgdaSymbol{=}\AgdaSpace{}%
\AgdaFunction{subst}\AgdaSpace{}%
\AgdaBound{P}\AgdaSpace{}%
\AgdaSymbol{(}\AgdaFunction{sym}\AgdaSpace{}%
\AgdaBound{x≡ℓ}\AgdaSymbol{)}\AgdaSpace{}%
\AgdaBound{pℓ}\<%
\\
\>[0]\AgdaSymbol{...}\AgdaSpace{}%
\AgdaSymbol{|}\AgdaSpace{}%
\AgdaInductiveConstructor{inj₂}\AgdaSpace{}%
\AgdaBound{x≡r}\AgdaSpace{}%
\AgdaSymbol{=}\AgdaSpace{}%
\AgdaFunction{subst}\AgdaSpace{}%
\AgdaBound{P}\AgdaSpace{}%
\AgdaSymbol{(}\AgdaFunction{sym}\AgdaSpace{}%
\AgdaBound{x≡r}\AgdaSymbol{)}\AgdaSpace{}%
\AgdaBound{pr}\<%
\end{code}

Surplus is therefore not only excluded but eliminated by construction. The
fourth field does the opposite work. \texttt{ℓ≠r} forbids the collapse of
the two positions into one. Anyone who tries to identify them produces a
contradiction the checker rejects.

Taken together, separation and cover yield a stable two. The next question
is what self-maps such a carrier admits.

% ─────────────────────────────────────────────────────────────────
\chapter{What Already Follows}
\label{ch:endomorphisms-closure}
% ─────────────────────────────────────────────────────────────────

Once the first distinction is given, the next question follows at once:
what self-maps does such a carrier admit? The point of this chapter is
modest in claim and strong in consequence. It does not yet deliver the
later closure reading of \textit{Void}. It shows that even at this first
formal threshold the space of endomorphisms is no longer open-ended.

Once the carrier is fixed to two classes, self-maps are almost entirely
determined by their values on \texttt{ℓ} and \texttt{r}. To compare two
maps, we use pointwise equality. On this carrier it is especially
manageable because \texttt{cover} reduces every argument to one of the two
boundary cases.

\begin{code}%
\>[0]\AgdaKeyword{infix}\AgdaSpace{}%
\AgdaNumber{4}\AgdaSpace{}%
\AgdaOperator{\AgdaFunction{\AgdaUnderscore{}≗\AgdaUnderscore{}}}\<%
\\
\>[0]\AgdaOperator{\AgdaFunction{\AgdaUnderscore{}≗\AgdaUnderscore{}}}\AgdaSpace{}%
\AgdaSymbol{:}\AgdaSpace{}%
\AgdaSymbol{\{}\AgdaBound{ℓ₁}\AgdaSpace{}%
\AgdaBound{ℓ₂}\AgdaSpace{}%
\AgdaSymbol{:}\AgdaSpace{}%
\AgdaPostulate{Level}\AgdaSymbol{\}}\AgdaSpace{}%
\AgdaSymbol{\{}\AgdaBound{A}\AgdaSpace{}%
\AgdaSymbol{:}\AgdaSpace{}%
\AgdaPrimitive{Set}\AgdaSpace{}%
\AgdaBound{ℓ₁}\AgdaSymbol{\}}\AgdaSpace{}%
\AgdaSymbol{\{}\AgdaBound{B}\AgdaSpace{}%
\AgdaSymbol{:}\AgdaSpace{}%
\AgdaPrimitive{Set}\AgdaSpace{}%
\AgdaBound{ℓ₂}\AgdaSymbol{\}}\AgdaSpace{}%
\AgdaSymbol{→}\AgdaSpace{}%
\AgdaSymbol{(}\AgdaBound{A}\AgdaSpace{}%
\AgdaSymbol{→}\AgdaSpace{}%
\AgdaBound{B}\AgdaSymbol{)}\AgdaSpace{}%
\AgdaSymbol{→}\AgdaSpace{}%
\AgdaSymbol{(}\AgdaBound{A}\AgdaSpace{}%
\AgdaSymbol{→}\AgdaSpace{}%
\AgdaBound{B}\AgdaSymbol{)}\AgdaSpace{}%
\AgdaSymbol{→}\AgdaSpace{}%
\AgdaPrimitive{Set}\AgdaSpace{}%
\AgdaSymbol{(}\AgdaBound{ℓ₁}\AgdaSpace{}%
\AgdaOperator{\AgdaPrimitive{⊔}}\AgdaSpace{}%
\AgdaBound{ℓ₂}\AgdaSymbol{)}\<%
\\
\>[0]\AgdaOperator{\AgdaFunction{\AgdaUnderscore{}≗\AgdaUnderscore{}}}\AgdaSpace{}%
\AgdaSymbol{\{}\AgdaArgument{A}\AgdaSpace{}%
\AgdaSymbol{=}\AgdaSpace{}%
\AgdaBound{A}\AgdaSymbol{\}}\AgdaSpace{}%
\AgdaBound{f}\AgdaSpace{}%
\AgdaBound{g}\AgdaSpace{}%
\AgdaSymbol{=}\AgdaSpace{}%
\AgdaSymbol{(}\AgdaBound{x}\AgdaSpace{}%
\AgdaSymbol{:}\AgdaSpace{}%
\AgdaBound{A}\AgdaSymbol{)}\AgdaSpace{}%
\AgdaSymbol{→}\AgdaSpace{}%
\AgdaBound{f}\AgdaSpace{}%
\AgdaBound{x}\AgdaSpace{}%
\AgdaOperator{\AgdaDatatype{≡}}\AgdaSpace{}%
\AgdaBound{g}\AgdaSpace{}%
\AgdaBound{x}\<%
\\
%
\\[\AgdaEmptyExtraSkip]%
\>[0]\AgdaFunction{id}\AgdaSpace{}%
\AgdaSymbol{:}\AgdaSpace{}%
\AgdaSymbol{\{}\AgdaBound{A}\AgdaSpace{}%
\AgdaSymbol{:}\AgdaSpace{}%
\AgdaPrimitive{Set}\AgdaSymbol{\}}\AgdaSpace{}%
\AgdaSymbol{→}\AgdaSpace{}%
\AgdaBound{A}\AgdaSpace{}%
\AgdaSymbol{→}\AgdaSpace{}%
\AgdaBound{A}\<%
\\
\>[0]\AgdaFunction{id}\AgdaSpace{}%
\AgdaBound{x}\AgdaSpace{}%
\AgdaSymbol{=}\AgdaSpace{}%
\AgdaBound{x}\<%
\end{code}

Among these maps one is especially important: the map that exchanges the
two positions, sending \texttt{ℓ} to \texttt{r} and \texttt{r} to
\texttt{ℓ}. Call it the dual. Applied twice, it returns every point to its
starting place.

\begin{code}%
\>[0]\AgdaFunction{Distinction-dual}\AgdaSpace{}%
\AgdaSymbol{:}\AgdaSpace{}%
\AgdaSymbol{(}\AgdaBound{d}\AgdaSpace{}%
\AgdaSymbol{:}\AgdaSpace{}%
\AgdaRecord{Distinction}\AgdaSymbol{)}\AgdaSpace{}%
\AgdaSymbol{→}\AgdaSpace{}%
\AgdaField{S}\AgdaSpace{}%
\AgdaBound{d}\AgdaSpace{}%
\AgdaSymbol{→}\AgdaSpace{}%
\AgdaField{S}\AgdaSpace{}%
\AgdaBound{d}\<%
\\
\>[0]\AgdaFunction{Distinction-dual}\AgdaSpace{}%
\AgdaBound{d}\AgdaSpace{}%
\AgdaBound{x}\AgdaSpace{}%
\AgdaKeyword{with}\AgdaSpace{}%
\AgdaField{cover}\AgdaSpace{}%
\AgdaBound{d}\AgdaSpace{}%
\AgdaBound{x}\<%
\\
\>[0]\AgdaSymbol{...}\AgdaSpace{}%
\AgdaSymbol{|}\AgdaSpace{}%
\AgdaInductiveConstructor{inj₁}\AgdaSpace{}%
\AgdaSymbol{\AgdaUnderscore{}}\AgdaSpace{}%
\AgdaSymbol{=}\AgdaSpace{}%
\AgdaField{r}\AgdaSpace{}%
\AgdaBound{d}\<%
\\
\>[0]\AgdaSymbol{...}\AgdaSpace{}%
\AgdaSymbol{|}\AgdaSpace{}%
\AgdaInductiveConstructor{inj₂}\AgdaSpace{}%
\AgdaSymbol{\AgdaUnderscore{}}\AgdaSpace{}%
\AgdaSymbol{=}\AgdaSpace{}%
\AgdaField{ℓ}\AgdaSpace{}%
\AgdaBound{d}\<%
\\
%
\\[\AgdaEmptyExtraSkip]%
\>[0]\AgdaFunction{Distinction-dual-involutive}\AgdaSpace{}%
\AgdaSymbol{:}\<%
\\
\>[0][@{}l@{\AgdaIndent{0}}]%
\>[2]\AgdaSymbol{(}\AgdaBound{d}\AgdaSpace{}%
\AgdaSymbol{:}\AgdaSpace{}%
\AgdaRecord{Distinction}\AgdaSymbol{)}\AgdaSpace{}%
\AgdaSymbol{→}\<%
\\
%
\>[2]\AgdaSymbol{(}\AgdaBound{x}\AgdaSpace{}%
\AgdaSymbol{:}\AgdaSpace{}%
\AgdaField{S}\AgdaSpace{}%
\AgdaBound{d}\AgdaSymbol{)}\AgdaSpace{}%
\AgdaSymbol{→}\<%
\\
%
\>[2]\AgdaFunction{Distinction-dual}\AgdaSpace{}%
\AgdaBound{d}\AgdaSpace{}%
\AgdaSymbol{(}\AgdaFunction{Distinction-dual}\AgdaSpace{}%
\AgdaBound{d}\AgdaSpace{}%
\AgdaBound{x}\AgdaSymbol{)}\AgdaSpace{}%
\AgdaOperator{\AgdaDatatype{≡}}\AgdaSpace{}%
\AgdaBound{x}\<%
\\
\>[0]\AgdaFunction{Distinction-dual-involutive}\AgdaSpace{}%
\AgdaBound{d}\AgdaSpace{}%
\AgdaSymbol{=}\<%
\\
\>[0][@{}l@{\AgdaIndent{0}}]%
\>[2]\AgdaFunction{Distinction-elim}\AgdaSpace{}%
\AgdaBound{d}\AgdaSpace{}%
\AgdaFunction{proof-ℓ}\AgdaSpace{}%
\AgdaFunction{proof-r}\<%
\\
%
\>[2]\AgdaKeyword{where}\<%
\\
\>[2][@{}l@{\AgdaIndent{0}}]%
\>[4]\AgdaFunction{proof-ℓ}\AgdaSpace{}%
\AgdaSymbol{:}\AgdaSpace{}%
\AgdaFunction{Distinction-dual}\AgdaSpace{}%
\AgdaBound{d}\AgdaSpace{}%
\AgdaSymbol{(}\AgdaFunction{Distinction-dual}\AgdaSpace{}%
\AgdaBound{d}\AgdaSpace{}%
\AgdaSymbol{(}\AgdaField{ℓ}\AgdaSpace{}%
\AgdaBound{d}\AgdaSymbol{))}\AgdaSpace{}%
\AgdaOperator{\AgdaDatatype{≡}}\AgdaSpace{}%
\AgdaField{ℓ}\AgdaSpace{}%
\AgdaBound{d}\<%
\\
%
\>[4]\AgdaFunction{proof-ℓ}\AgdaSpace{}%
\AgdaKeyword{with}\AgdaSpace{}%
\AgdaField{cover}\AgdaSpace{}%
\AgdaBound{d}\AgdaSpace{}%
\AgdaSymbol{(}\AgdaField{ℓ}\AgdaSpace{}%
\AgdaBound{d}\AgdaSymbol{)}\<%
\\
%
\>[4]\AgdaFunction{...}\AgdaSpace{}%
\AgdaSymbol{|}\AgdaSpace{}%
\AgdaInductiveConstructor{inj₂}\AgdaSpace{}%
\AgdaBound{ℓ≡r}\AgdaSpace{}%
\AgdaSymbol{=}\AgdaSpace{}%
\AgdaFunction{⊥-elim}\AgdaSpace{}%
\AgdaSymbol{((}\AgdaField{ℓ≠r}\AgdaSpace{}%
\AgdaBound{d}\AgdaSymbol{)}\AgdaSpace{}%
\AgdaBound{ℓ≡r}\AgdaSymbol{)}\<%
\\
%
\>[4]\AgdaFunction{...}\AgdaSpace{}%
\AgdaSymbol{|}\AgdaSpace{}%
\AgdaInductiveConstructor{inj₁}\AgdaSpace{}%
\AgdaSymbol{\AgdaUnderscore{}}\AgdaSpace{}%
\AgdaKeyword{with}\AgdaSpace{}%
\AgdaField{cover}\AgdaSpace{}%
\AgdaBound{d}\AgdaSpace{}%
\AgdaSymbol{(}\AgdaField{r}\AgdaSpace{}%
\AgdaBound{d}\AgdaSymbol{)}\<%
\\
%
\>[4]\AgdaSymbol{...}\AgdaSpace{}%
\AgdaSymbol{|}\AgdaSpace{}%
\AgdaInductiveConstructor{inj₁}\AgdaSpace{}%
\AgdaBound{r≡ℓ}\AgdaSpace{}%
\AgdaSymbol{=}\AgdaSpace{}%
\AgdaFunction{⊥-elim}\AgdaSpace{}%
\AgdaSymbol{((}\AgdaField{ℓ≠r}\AgdaSpace{}%
\AgdaBound{d}\AgdaSymbol{)}\AgdaSpace{}%
\AgdaSymbol{(}\AgdaFunction{sym}\AgdaSpace{}%
\AgdaBound{r≡ℓ}\AgdaSymbol{))}\<%
\\
%
\>[4]\AgdaSymbol{...}\AgdaSpace{}%
\AgdaSymbol{|}\AgdaSpace{}%
\AgdaInductiveConstructor{inj₂}\AgdaSpace{}%
\AgdaSymbol{\AgdaUnderscore{}}%
\>[19]\AgdaSymbol{=}\AgdaSpace{}%
\AgdaInductiveConstructor{refl}\<%
\\
%
\\[\AgdaEmptyExtraSkip]%
%
\>[4]\AgdaFunction{proof-r}\AgdaSpace{}%
\AgdaSymbol{:}\AgdaSpace{}%
\AgdaFunction{Distinction-dual}\AgdaSpace{}%
\AgdaBound{d}\AgdaSpace{}%
\AgdaSymbol{(}\AgdaFunction{Distinction-dual}\AgdaSpace{}%
\AgdaBound{d}\AgdaSpace{}%
\AgdaSymbol{(}\AgdaField{r}\AgdaSpace{}%
\AgdaBound{d}\AgdaSymbol{))}\AgdaSpace{}%
\AgdaOperator{\AgdaDatatype{≡}}\AgdaSpace{}%
\AgdaField{r}\AgdaSpace{}%
\AgdaBound{d}\<%
\\
%
\>[4]\AgdaFunction{proof-r}\AgdaSpace{}%
\AgdaKeyword{with}\AgdaSpace{}%
\AgdaField{cover}\AgdaSpace{}%
\AgdaBound{d}\AgdaSpace{}%
\AgdaSymbol{(}\AgdaField{r}\AgdaSpace{}%
\AgdaBound{d}\AgdaSymbol{)}\<%
\\
%
\>[4]\AgdaFunction{...}\AgdaSpace{}%
\AgdaSymbol{|}\AgdaSpace{}%
\AgdaInductiveConstructor{inj₁}\AgdaSpace{}%
\AgdaBound{r≡ℓ}\AgdaSpace{}%
\AgdaSymbol{=}\AgdaSpace{}%
\AgdaFunction{⊥-elim}\AgdaSpace{}%
\AgdaSymbol{((}\AgdaField{ℓ≠r}\AgdaSpace{}%
\AgdaBound{d}\AgdaSymbol{)}\AgdaSpace{}%
\AgdaSymbol{(}\AgdaFunction{sym}\AgdaSpace{}%
\AgdaBound{r≡ℓ}\AgdaSymbol{))}\<%
\\
%
\>[4]\AgdaFunction{...}\AgdaSpace{}%
\AgdaSymbol{|}\AgdaSpace{}%
\AgdaInductiveConstructor{inj₂}\AgdaSpace{}%
\AgdaSymbol{\AgdaUnderscore{}}\AgdaSpace{}%
\AgdaKeyword{with}\AgdaSpace{}%
\AgdaField{cover}\AgdaSpace{}%
\AgdaBound{d}\AgdaSpace{}%
\AgdaSymbol{(}\AgdaField{ℓ}\AgdaSpace{}%
\AgdaBound{d}\AgdaSymbol{)}\<%
\\
%
\>[4]\AgdaSymbol{...}\AgdaSpace{}%
\AgdaSymbol{|}\AgdaSpace{}%
\AgdaInductiveConstructor{inj₂}\AgdaSpace{}%
\AgdaBound{ℓ≡r}\AgdaSpace{}%
\AgdaSymbol{=}\AgdaSpace{}%
\AgdaFunction{⊥-elim}\AgdaSpace{}%
\AgdaSymbol{((}\AgdaField{ℓ≠r}\AgdaSpace{}%
\AgdaBound{d}\AgdaSymbol{)}\AgdaSpace{}%
\AgdaBound{ℓ≡r}\AgdaSymbol{)}\<%
\\
%
\>[4]\AgdaSymbol{...}\AgdaSpace{}%
\AgdaSymbol{|}\AgdaSpace{}%
\AgdaInductiveConstructor{inj₁}\AgdaSpace{}%
\AgdaSymbol{\AgdaUnderscore{}}%
\>[19]\AgdaSymbol{=}\AgdaSpace{}%
\AgdaInductiveConstructor{refl}\<%
\end{code}

The central question is: how many maps are there, counted up to pointwise
equality? The boundary table allows exactly four combinations, one for each
pair of choices at \texttt{ℓ} and \texttt{r} --- the constant-left map,
the constant-right map, the identity, and the dual. We name those four
cases explicitly.

\begin{code}%
\>[0]\AgdaKeyword{data}\AgdaSpace{}%
\AgdaDatatype{EndoCase}\AgdaSpace{}%
\AgdaSymbol{:}\AgdaSpace{}%
\AgdaPrimitive{Set}\AgdaSpace{}%
\AgdaKeyword{where}\<%
\\
\>[0][@{}l@{\AgdaIndent{0}}]%
\>[2]\AgdaInductiveConstructor{case-constL}\AgdaSpace{}%
\AgdaSymbol{:}\AgdaSpace{}%
\AgdaDatatype{EndoCase}\<%
\\
%
\>[2]\AgdaInductiveConstructor{case-constR}\AgdaSpace{}%
\AgdaSymbol{:}\AgdaSpace{}%
\AgdaDatatype{EndoCase}\<%
\\
%
\>[2]\AgdaInductiveConstructor{case-id}%
\>[14]\AgdaSymbol{:}\AgdaSpace{}%
\AgdaDatatype{EndoCase}\<%
\\
%
\>[2]\AgdaInductiveConstructor{case-dual}%
\>[14]\AgdaSymbol{:}\AgdaSpace{}%
\AgdaDatatype{EndoCase}\<%
\end{code}

The four names must now be tied to actual maps. Inside a module
parametrized by a fixed distinction, \texttt{interpret} turns each label
into its endomorphism, and \texttt{classify} reads the matching label back
from an endomorphism's two boundary values.

\begin{code}%
\>[0]\AgdaKeyword{module}\AgdaSpace{}%
\AgdaModule{K₄}\AgdaSpace{}%
\AgdaSymbol{(}\AgdaBound{d}\AgdaSpace{}%
\AgdaSymbol{:}\AgdaSpace{}%
\AgdaRecord{Distinction}\AgdaSymbol{)}\AgdaSpace{}%
\AgdaKeyword{where}\<%
\\
\>[0][@{}l@{\AgdaIndent{0}}]%
\>[2]\AgdaFunction{Endo}\AgdaSpace{}%
\AgdaSymbol{:}\AgdaSpace{}%
\AgdaPrimitive{Set}\<%
\\
%
\>[2]\AgdaFunction{Endo}\AgdaSpace{}%
\AgdaSymbol{=}\AgdaSpace{}%
\AgdaField{S}\AgdaSpace{}%
\AgdaBound{d}\AgdaSpace{}%
\AgdaSymbol{→}\AgdaSpace{}%
\AgdaField{S}\AgdaSpace{}%
\AgdaBound{d}\<%
\\
%
\\[\AgdaEmptyExtraSkip]%
%
\>[2]\AgdaFunction{≗-refl}\AgdaSpace{}%
\AgdaSymbol{:}\AgdaSpace{}%
\AgdaSymbol{\{}\AgdaBound{f}\AgdaSpace{}%
\AgdaSymbol{:}\AgdaSpace{}%
\AgdaFunction{Endo}\AgdaSymbol{\}}\AgdaSpace{}%
\AgdaSymbol{→}\AgdaSpace{}%
\AgdaBound{f}\AgdaSpace{}%
\AgdaOperator{\AgdaFunction{≗}}\AgdaSpace{}%
\AgdaBound{f}\<%
\\
%
\>[2]\AgdaFunction{≗-refl}\AgdaSpace{}%
\AgdaBound{x}\AgdaSpace{}%
\AgdaSymbol{=}\AgdaSpace{}%
\AgdaInductiveConstructor{refl}\<%
\\
%
\\[\AgdaEmptyExtraSkip]%
%
\>[2]\AgdaFunction{≗-sym}\AgdaSpace{}%
\AgdaSymbol{:}\AgdaSpace{}%
\AgdaSymbol{\{}\AgdaBound{f}\AgdaSpace{}%
\AgdaBound{g}\AgdaSpace{}%
\AgdaSymbol{:}\AgdaSpace{}%
\AgdaFunction{Endo}\AgdaSymbol{\}}\AgdaSpace{}%
\AgdaSymbol{→}\AgdaSpace{}%
\AgdaBound{f}\AgdaSpace{}%
\AgdaOperator{\AgdaFunction{≗}}\AgdaSpace{}%
\AgdaBound{g}\AgdaSpace{}%
\AgdaSymbol{→}\AgdaSpace{}%
\AgdaBound{g}\AgdaSpace{}%
\AgdaOperator{\AgdaFunction{≗}}\AgdaSpace{}%
\AgdaBound{f}\<%
\\
%
\>[2]\AgdaFunction{≗-sym}\AgdaSpace{}%
\AgdaBound{p}\AgdaSpace{}%
\AgdaBound{x}\AgdaSpace{}%
\AgdaSymbol{=}\AgdaSpace{}%
\AgdaFunction{sym}\AgdaSpace{}%
\AgdaSymbol{(}\AgdaBound{p}\AgdaSpace{}%
\AgdaBound{x}\AgdaSymbol{)}\<%
\\
%
\\[\AgdaEmptyExtraSkip]%
%
\>[2]\AgdaFunction{≗-trans}\AgdaSpace{}%
\AgdaSymbol{:}\AgdaSpace{}%
\AgdaSymbol{\{}\AgdaBound{f}\AgdaSpace{}%
\AgdaBound{g}\AgdaSpace{}%
\AgdaBound{h}\AgdaSpace{}%
\AgdaSymbol{:}\AgdaSpace{}%
\AgdaFunction{Endo}\AgdaSymbol{\}}\AgdaSpace{}%
\AgdaSymbol{→}\AgdaSpace{}%
\AgdaBound{f}\AgdaSpace{}%
\AgdaOperator{\AgdaFunction{≗}}\AgdaSpace{}%
\AgdaBound{g}\AgdaSpace{}%
\AgdaSymbol{→}\AgdaSpace{}%
\AgdaBound{g}\AgdaSpace{}%
\AgdaOperator{\AgdaFunction{≗}}\AgdaSpace{}%
\AgdaBound{h}\AgdaSpace{}%
\AgdaSymbol{→}\AgdaSpace{}%
\AgdaBound{f}\AgdaSpace{}%
\AgdaOperator{\AgdaFunction{≗}}\AgdaSpace{}%
\AgdaBound{h}\<%
\\
%
\>[2]\AgdaFunction{≗-trans}\AgdaSpace{}%
\AgdaBound{p}\AgdaSpace{}%
\AgdaBound{q}\AgdaSpace{}%
\AgdaBound{x}\AgdaSpace{}%
\AgdaSymbol{=}\AgdaSpace{}%
\AgdaFunction{trans}\AgdaSpace{}%
\AgdaSymbol{(}\AgdaBound{p}\AgdaSpace{}%
\AgdaBound{x}\AgdaSymbol{)}\AgdaSpace{}%
\AgdaSymbol{(}\AgdaBound{q}\AgdaSpace{}%
\AgdaBound{x}\AgdaSymbol{)}\<%
\\
%
\\[\AgdaEmptyExtraSkip]%
%
\>[2]\AgdaFunction{constL}\AgdaSpace{}%
\AgdaSymbol{:}\AgdaSpace{}%
\AgdaFunction{Endo}\<%
\\
%
\>[2]\AgdaFunction{constL}\AgdaSpace{}%
\AgdaSymbol{\AgdaUnderscore{}}\AgdaSpace{}%
\AgdaSymbol{=}\AgdaSpace{}%
\AgdaField{ℓ}\AgdaSpace{}%
\AgdaBound{d}\<%
\\
%
\\[\AgdaEmptyExtraSkip]%
%
\>[2]\AgdaFunction{constR}\AgdaSpace{}%
\AgdaSymbol{:}\AgdaSpace{}%
\AgdaFunction{Endo}\<%
\\
%
\>[2]\AgdaFunction{constR}\AgdaSpace{}%
\AgdaSymbol{\AgdaUnderscore{}}\AgdaSpace{}%
\AgdaSymbol{=}\AgdaSpace{}%
\AgdaField{r}\AgdaSpace{}%
\AgdaBound{d}\<%
\\
%
\\[\AgdaEmptyExtraSkip]%
%
\>[2]\AgdaFunction{dual}\AgdaSpace{}%
\AgdaSymbol{:}\AgdaSpace{}%
\AgdaFunction{Endo}\<%
\\
%
\>[2]\AgdaFunction{dual}\AgdaSpace{}%
\AgdaSymbol{=}\AgdaSpace{}%
\AgdaFunction{Distinction-dual}\AgdaSpace{}%
\AgdaBound{d}\<%
\\
%
\\[\AgdaEmptyExtraSkip]%
%
\>[2]\AgdaFunction{dual-ℓ}\AgdaSpace{}%
\AgdaSymbol{:}\AgdaSpace{}%
\AgdaFunction{dual}\AgdaSpace{}%
\AgdaSymbol{(}\AgdaField{ℓ}\AgdaSpace{}%
\AgdaBound{d}\AgdaSymbol{)}\AgdaSpace{}%
\AgdaOperator{\AgdaDatatype{≡}}\AgdaSpace{}%
\AgdaField{r}\AgdaSpace{}%
\AgdaBound{d}\<%
\\
%
\>[2]\AgdaFunction{dual-ℓ}\AgdaSpace{}%
\AgdaKeyword{with}\AgdaSpace{}%
\AgdaField{cover}\AgdaSpace{}%
\AgdaBound{d}\AgdaSpace{}%
\AgdaSymbol{(}\AgdaField{ℓ}\AgdaSpace{}%
\AgdaBound{d}\AgdaSymbol{)}\<%
\\
%
\>[2]\AgdaFunction{...}\AgdaSpace{}%
\AgdaSymbol{|}\AgdaSpace{}%
\AgdaInductiveConstructor{inj₁}\AgdaSpace{}%
\AgdaSymbol{\AgdaUnderscore{}}%
\>[17]\AgdaSymbol{=}\AgdaSpace{}%
\AgdaInductiveConstructor{refl}\<%
\\
%
\>[2]\AgdaFunction{...}\AgdaSpace{}%
\AgdaSymbol{|}\AgdaSpace{}%
\AgdaInductiveConstructor{inj₂}\AgdaSpace{}%
\AgdaBound{ℓ≡r}\AgdaSpace{}%
\AgdaSymbol{=}\AgdaSpace{}%
\AgdaFunction{⊥-elim}\AgdaSpace{}%
\AgdaSymbol{((}\AgdaField{ℓ≠r}\AgdaSpace{}%
\AgdaBound{d}\AgdaSymbol{)}\AgdaSpace{}%
\AgdaBound{ℓ≡r}\AgdaSymbol{)}\<%
\\
%
\\[\AgdaEmptyExtraSkip]%
%
\>[2]\AgdaFunction{dual-r}\AgdaSpace{}%
\AgdaSymbol{:}\AgdaSpace{}%
\AgdaFunction{dual}\AgdaSpace{}%
\AgdaSymbol{(}\AgdaField{r}\AgdaSpace{}%
\AgdaBound{d}\AgdaSymbol{)}\AgdaSpace{}%
\AgdaOperator{\AgdaDatatype{≡}}\AgdaSpace{}%
\AgdaField{ℓ}\AgdaSpace{}%
\AgdaBound{d}\<%
\\
%
\>[2]\AgdaFunction{dual-r}\AgdaSpace{}%
\AgdaKeyword{with}\AgdaSpace{}%
\AgdaField{cover}\AgdaSpace{}%
\AgdaBound{d}\AgdaSpace{}%
\AgdaSymbol{(}\AgdaField{r}\AgdaSpace{}%
\AgdaBound{d}\AgdaSymbol{)}\<%
\\
%
\>[2]\AgdaFunction{...}\AgdaSpace{}%
\AgdaSymbol{|}\AgdaSpace{}%
\AgdaInductiveConstructor{inj₁}\AgdaSpace{}%
\AgdaBound{r≡ℓ}\AgdaSpace{}%
\AgdaSymbol{=}\AgdaSpace{}%
\AgdaFunction{⊥-elim}\AgdaSpace{}%
\AgdaSymbol{((}\AgdaField{ℓ≠r}\AgdaSpace{}%
\AgdaBound{d}\AgdaSymbol{)}\AgdaSpace{}%
\AgdaSymbol{(}\AgdaFunction{sym}\AgdaSpace{}%
\AgdaBound{r≡ℓ}\AgdaSymbol{))}\<%
\\
%
\>[2]\AgdaFunction{...}\AgdaSpace{}%
\AgdaSymbol{|}\AgdaSpace{}%
\AgdaInductiveConstructor{inj₂}\AgdaSpace{}%
\AgdaSymbol{\AgdaUnderscore{}}%
\>[17]\AgdaSymbol{=}\AgdaSpace{}%
\AgdaInductiveConstructor{refl}\<%
\\
%
\\[\AgdaEmptyExtraSkip]%
%
\>[2]\AgdaFunction{interpret}\AgdaSpace{}%
\AgdaSymbol{:}\AgdaSpace{}%
\AgdaDatatype{EndoCase}\AgdaSpace{}%
\AgdaSymbol{→}\AgdaSpace{}%
\AgdaFunction{Endo}\<%
\\
%
\>[2]\AgdaFunction{interpret}\AgdaSpace{}%
\AgdaInductiveConstructor{case-constL}\AgdaSpace{}%
\AgdaSymbol{=}\AgdaSpace{}%
\AgdaFunction{constL}\<%
\\
%
\>[2]\AgdaFunction{interpret}\AgdaSpace{}%
\AgdaInductiveConstructor{case-constR}\AgdaSpace{}%
\AgdaSymbol{=}\AgdaSpace{}%
\AgdaFunction{constR}\<%
\\
%
\>[2]\AgdaFunction{interpret}\AgdaSpace{}%
\AgdaInductiveConstructor{case-id}%
\>[24]\AgdaSymbol{=}\AgdaSpace{}%
\AgdaFunction{id}\<%
\\
%
\>[2]\AgdaFunction{interpret}\AgdaSpace{}%
\AgdaInductiveConstructor{case-dual}%
\>[24]\AgdaSymbol{=}\AgdaSpace{}%
\AgdaFunction{dual}\<%
\\
%
\\[\AgdaEmptyExtraSkip]%
%
\>[2]\AgdaFunction{classify}\AgdaSpace{}%
\AgdaSymbol{:}\AgdaSpace{}%
\AgdaFunction{Endo}\AgdaSpace{}%
\AgdaSymbol{→}\AgdaSpace{}%
\AgdaDatatype{EndoCase}\<%
\\
%
\>[2]\AgdaFunction{classify}\AgdaSpace{}%
\AgdaBound{f}\AgdaSpace{}%
\AgdaKeyword{with}\AgdaSpace{}%
\AgdaField{cover}\AgdaSpace{}%
\AgdaBound{d}\AgdaSpace{}%
\AgdaSymbol{(}\AgdaBound{f}\AgdaSpace{}%
\AgdaSymbol{(}\AgdaField{ℓ}\AgdaSpace{}%
\AgdaBound{d}\AgdaSymbol{))}\AgdaSpace{}%
\AgdaSymbol{|}\AgdaSpace{}%
\AgdaField{cover}\AgdaSpace{}%
\AgdaBound{d}\AgdaSpace{}%
\AgdaSymbol{(}\AgdaBound{f}\AgdaSpace{}%
\AgdaSymbol{(}\AgdaField{r}\AgdaSpace{}%
\AgdaBound{d}\AgdaSymbol{))}\<%
\\
%
\>[2]\AgdaSymbol{...}\AgdaSpace{}%
\AgdaSymbol{|}\AgdaSpace{}%
\AgdaInductiveConstructor{inj₁}\AgdaSpace{}%
\AgdaSymbol{\AgdaUnderscore{}}\AgdaSpace{}%
\AgdaSymbol{|}\AgdaSpace{}%
\AgdaInductiveConstructor{inj₁}\AgdaSpace{}%
\AgdaSymbol{\AgdaUnderscore{}}\AgdaSpace{}%
\AgdaSymbol{=}\AgdaSpace{}%
\AgdaInductiveConstructor{case-constL}\<%
\\
%
\>[2]\AgdaSymbol{...}\AgdaSpace{}%
\AgdaSymbol{|}\AgdaSpace{}%
\AgdaInductiveConstructor{inj₂}\AgdaSpace{}%
\AgdaSymbol{\AgdaUnderscore{}}\AgdaSpace{}%
\AgdaSymbol{|}\AgdaSpace{}%
\AgdaInductiveConstructor{inj₂}\AgdaSpace{}%
\AgdaSymbol{\AgdaUnderscore{}}\AgdaSpace{}%
\AgdaSymbol{=}\AgdaSpace{}%
\AgdaInductiveConstructor{case-constR}\<%
\\
%
\>[2]\AgdaSymbol{...}\AgdaSpace{}%
\AgdaSymbol{|}\AgdaSpace{}%
\AgdaInductiveConstructor{inj₁}\AgdaSpace{}%
\AgdaSymbol{\AgdaUnderscore{}}\AgdaSpace{}%
\AgdaSymbol{|}\AgdaSpace{}%
\AgdaInductiveConstructor{inj₂}\AgdaSpace{}%
\AgdaSymbol{\AgdaUnderscore{}}\AgdaSpace{}%
\AgdaSymbol{=}\AgdaSpace{}%
\AgdaInductiveConstructor{case-id}\<%
\\
%
\>[2]\AgdaSymbol{...}\AgdaSpace{}%
\AgdaSymbol{|}\AgdaSpace{}%
\AgdaInductiveConstructor{inj₂}\AgdaSpace{}%
\AgdaSymbol{\AgdaUnderscore{}}\AgdaSpace{}%
\AgdaSymbol{|}\AgdaSpace{}%
\AgdaInductiveConstructor{inj₁}\AgdaSpace{}%
\AgdaSymbol{\AgdaUnderscore{}}\AgdaSpace{}%
\AgdaSymbol{=}\AgdaSpace{}%
\AgdaInductiveConstructor{case-dual}\<%
\end{code}

Reading a function's label and then interpreting it returns the same
function: at the two boundary points by direct inspection, and therefore
everywhere by the eliminator. The two boundary lemmas this rests on are
routine case analyses; they are part of the checked file and are not
reprinted here.

\begin{code}[hide]%
%
\>[2]\AgdaFunction{sound-at-ℓ}\AgdaSpace{}%
\AgdaSymbol{:}\AgdaSpace{}%
\AgdaSymbol{(}\AgdaBound{f}\AgdaSpace{}%
\AgdaSymbol{:}\AgdaSpace{}%
\AgdaFunction{Endo}\AgdaSymbol{)}\AgdaSpace{}%
\AgdaSymbol{→}\AgdaSpace{}%
\AgdaFunction{interpret}\AgdaSpace{}%
\AgdaSymbol{(}\AgdaFunction{classify}\AgdaSpace{}%
\AgdaBound{f}\AgdaSymbol{)}\AgdaSpace{}%
\AgdaSymbol{(}\AgdaField{ℓ}\AgdaSpace{}%
\AgdaBound{d}\AgdaSymbol{)}\AgdaSpace{}%
\AgdaOperator{\AgdaDatatype{≡}}\AgdaSpace{}%
\AgdaBound{f}\AgdaSpace{}%
\AgdaSymbol{(}\AgdaField{ℓ}\AgdaSpace{}%
\AgdaBound{d}\AgdaSymbol{)}\<%
\\
%
\>[2]\AgdaFunction{sound-at-ℓ}\AgdaSpace{}%
\AgdaBound{f}\AgdaSpace{}%
\AgdaKeyword{with}\AgdaSpace{}%
\AgdaField{cover}\AgdaSpace{}%
\AgdaBound{d}\AgdaSpace{}%
\AgdaSymbol{(}\AgdaBound{f}\AgdaSpace{}%
\AgdaSymbol{(}\AgdaField{ℓ}\AgdaSpace{}%
\AgdaBound{d}\AgdaSymbol{))}\AgdaSpace{}%
\AgdaSymbol{|}\AgdaSpace{}%
\AgdaField{cover}\AgdaSpace{}%
\AgdaBound{d}\AgdaSpace{}%
\AgdaSymbol{(}\AgdaBound{f}\AgdaSpace{}%
\AgdaSymbol{(}\AgdaField{r}\AgdaSpace{}%
\AgdaBound{d}\AgdaSymbol{))}\<%
\\
%
\>[2]\AgdaSymbol{...}\AgdaSpace{}%
\AgdaSymbol{|}\AgdaSpace{}%
\AgdaInductiveConstructor{inj₁}\AgdaSpace{}%
\AgdaBound{fl≡ℓ}\AgdaSpace{}%
\AgdaSymbol{|}\AgdaSpace{}%
\AgdaInductiveConstructor{inj₁}\AgdaSpace{}%
\AgdaSymbol{\AgdaUnderscore{}}%
\>[31]\AgdaSymbol{=}\AgdaSpace{}%
\AgdaFunction{sym}\AgdaSpace{}%
\AgdaBound{fl≡ℓ}\<%
\\
%
\>[2]\AgdaSymbol{...}\AgdaSpace{}%
\AgdaSymbol{|}\AgdaSpace{}%
\AgdaInductiveConstructor{inj₂}\AgdaSpace{}%
\AgdaBound{fl≡r}\AgdaSpace{}%
\AgdaSymbol{|}\AgdaSpace{}%
\AgdaInductiveConstructor{inj₂}\AgdaSpace{}%
\AgdaSymbol{\AgdaUnderscore{}}%
\>[31]\AgdaSymbol{=}\AgdaSpace{}%
\AgdaFunction{sym}\AgdaSpace{}%
\AgdaBound{fl≡r}\<%
\\
%
\>[2]\AgdaSymbol{...}\AgdaSpace{}%
\AgdaSymbol{|}\AgdaSpace{}%
\AgdaInductiveConstructor{inj₁}\AgdaSpace{}%
\AgdaBound{fl≡ℓ}\AgdaSpace{}%
\AgdaSymbol{|}\AgdaSpace{}%
\AgdaInductiveConstructor{inj₂}\AgdaSpace{}%
\AgdaSymbol{\AgdaUnderscore{}}%
\>[31]\AgdaSymbol{=}\AgdaSpace{}%
\AgdaFunction{sym}\AgdaSpace{}%
\AgdaBound{fl≡ℓ}\<%
\\
%
\>[2]\AgdaSymbol{...}\AgdaSpace{}%
\AgdaSymbol{|}\AgdaSpace{}%
\AgdaInductiveConstructor{inj₂}\AgdaSpace{}%
\AgdaBound{fl≡r}\AgdaSpace{}%
\AgdaSymbol{|}\AgdaSpace{}%
\AgdaInductiveConstructor{inj₁}\AgdaSpace{}%
\AgdaSymbol{\AgdaUnderscore{}}%
\>[31]\AgdaSymbol{=}\AgdaSpace{}%
\AgdaFunction{trans}\AgdaSpace{}%
\AgdaFunction{dual-ℓ}\AgdaSpace{}%
\AgdaSymbol{(}\AgdaFunction{sym}\AgdaSpace{}%
\AgdaBound{fl≡r}\AgdaSymbol{)}\<%
\\
%
\\[\AgdaEmptyExtraSkip]%
%
\>[2]\AgdaFunction{sound-at-r}\AgdaSpace{}%
\AgdaSymbol{:}\AgdaSpace{}%
\AgdaSymbol{(}\AgdaBound{f}\AgdaSpace{}%
\AgdaSymbol{:}\AgdaSpace{}%
\AgdaFunction{Endo}\AgdaSymbol{)}\AgdaSpace{}%
\AgdaSymbol{→}\AgdaSpace{}%
\AgdaFunction{interpret}\AgdaSpace{}%
\AgdaSymbol{(}\AgdaFunction{classify}\AgdaSpace{}%
\AgdaBound{f}\AgdaSymbol{)}\AgdaSpace{}%
\AgdaSymbol{(}\AgdaField{r}\AgdaSpace{}%
\AgdaBound{d}\AgdaSymbol{)}\AgdaSpace{}%
\AgdaOperator{\AgdaDatatype{≡}}\AgdaSpace{}%
\AgdaBound{f}\AgdaSpace{}%
\AgdaSymbol{(}\AgdaField{r}\AgdaSpace{}%
\AgdaBound{d}\AgdaSymbol{)}\<%
\\
%
\>[2]\AgdaFunction{sound-at-r}\AgdaSpace{}%
\AgdaBound{f}\AgdaSpace{}%
\AgdaKeyword{with}\AgdaSpace{}%
\AgdaField{cover}\AgdaSpace{}%
\AgdaBound{d}\AgdaSpace{}%
\AgdaSymbol{(}\AgdaBound{f}\AgdaSpace{}%
\AgdaSymbol{(}\AgdaField{ℓ}\AgdaSpace{}%
\AgdaBound{d}\AgdaSymbol{))}\AgdaSpace{}%
\AgdaSymbol{|}\AgdaSpace{}%
\AgdaField{cover}\AgdaSpace{}%
\AgdaBound{d}\AgdaSpace{}%
\AgdaSymbol{(}\AgdaBound{f}\AgdaSpace{}%
\AgdaSymbol{(}\AgdaField{r}\AgdaSpace{}%
\AgdaBound{d}\AgdaSymbol{))}\<%
\\
%
\>[2]\AgdaSymbol{...}\AgdaSpace{}%
\AgdaSymbol{|}\AgdaSpace{}%
\AgdaInductiveConstructor{inj₁}\AgdaSpace{}%
\AgdaSymbol{\AgdaUnderscore{}}%
\>[19]\AgdaSymbol{|}\AgdaSpace{}%
\AgdaInductiveConstructor{inj₁}\AgdaSpace{}%
\AgdaBound{fr≡ℓ}\AgdaSpace{}%
\AgdaSymbol{=}\AgdaSpace{}%
\AgdaFunction{sym}\AgdaSpace{}%
\AgdaBound{fr≡ℓ}\<%
\\
%
\>[2]\AgdaSymbol{...}\AgdaSpace{}%
\AgdaSymbol{|}\AgdaSpace{}%
\AgdaInductiveConstructor{inj₂}\AgdaSpace{}%
\AgdaSymbol{\AgdaUnderscore{}}%
\>[19]\AgdaSymbol{|}\AgdaSpace{}%
\AgdaInductiveConstructor{inj₂}\AgdaSpace{}%
\AgdaBound{fr≡r}\AgdaSpace{}%
\AgdaSymbol{=}\AgdaSpace{}%
\AgdaFunction{sym}\AgdaSpace{}%
\AgdaBound{fr≡r}\<%
\\
%
\>[2]\AgdaSymbol{...}\AgdaSpace{}%
\AgdaSymbol{|}\AgdaSpace{}%
\AgdaInductiveConstructor{inj₁}\AgdaSpace{}%
\AgdaSymbol{\AgdaUnderscore{}}%
\>[19]\AgdaSymbol{|}\AgdaSpace{}%
\AgdaInductiveConstructor{inj₂}\AgdaSpace{}%
\AgdaBound{fr≡r}\AgdaSpace{}%
\AgdaSymbol{=}\AgdaSpace{}%
\AgdaFunction{sym}\AgdaSpace{}%
\AgdaBound{fr≡r}\<%
\\
%
\>[2]\AgdaSymbol{...}\AgdaSpace{}%
\AgdaSymbol{|}\AgdaSpace{}%
\AgdaInductiveConstructor{inj₂}\AgdaSpace{}%
\AgdaSymbol{\AgdaUnderscore{}}%
\>[19]\AgdaSymbol{|}\AgdaSpace{}%
\AgdaInductiveConstructor{inj₁}\AgdaSpace{}%
\AgdaBound{fr≡ℓ}\AgdaSpace{}%
\AgdaSymbol{=}\AgdaSpace{}%
\AgdaFunction{trans}\AgdaSpace{}%
\AgdaFunction{dual-r}\AgdaSpace{}%
\AgdaSymbol{(}\AgdaFunction{sym}\AgdaSpace{}%
\AgdaBound{fr≡ℓ}\AgdaSymbol{)}\<%
\end{code}

\begin{code}%
%
\>[2]\AgdaFunction{classify-sound}\AgdaSpace{}%
\AgdaSymbol{:}\AgdaSpace{}%
\AgdaSymbol{(}\AgdaBound{f}\AgdaSpace{}%
\AgdaSymbol{:}\AgdaSpace{}%
\AgdaFunction{Endo}\AgdaSymbol{)}\AgdaSpace{}%
\AgdaSymbol{→}\AgdaSpace{}%
\AgdaFunction{interpret}\AgdaSpace{}%
\AgdaSymbol{(}\AgdaFunction{classify}\AgdaSpace{}%
\AgdaBound{f}\AgdaSymbol{)}\AgdaSpace{}%
\AgdaOperator{\AgdaFunction{≗}}\AgdaSpace{}%
\AgdaBound{f}\<%
\\
%
\>[2]\AgdaFunction{classify-sound}\AgdaSpace{}%
\AgdaBound{f}\AgdaSpace{}%
\AgdaBound{x}\AgdaSpace{}%
\AgdaSymbol{=}\AgdaSpace{}%
\AgdaFunction{Distinction-elim}\AgdaSpace{}%
\AgdaBound{d}\AgdaSpace{}%
\AgdaSymbol{(}\AgdaFunction{sound-at-ℓ}\AgdaSpace{}%
\AgdaBound{f}\AgdaSymbol{)}\AgdaSpace{}%
\AgdaSymbol{(}\AgdaFunction{sound-at-r}\AgdaSpace{}%
\AgdaBound{f}\AgdaSymbol{)}\AgdaSpace{}%
\AgdaBound{x}\<%
\end{code}

The label is not merely available; it is forced. No two distinct labels
can interpret to the same function --- sixteen comparisons, one for each
ordered pair of labels, every one discharged by the single separation
witness \texttt{ℓ≠r}. That case analysis is routine and lives in the
checked file without crowding the page. From it, classification is
unique, and reading a label out and back in returns the very same label.

\begin{code}[hide]%
%
\>[2]\AgdaFunction{interpret-injective}\AgdaSpace{}%
\AgdaSymbol{:}\<%
\\
\>[2][@{}l@{\AgdaIndent{0}}]%
\>[4]\AgdaSymbol{(}\AgdaBound{c}\AgdaSpace{}%
\AgdaBound{c'}\AgdaSpace{}%
\AgdaSymbol{:}\AgdaSpace{}%
\AgdaDatatype{EndoCase}\AgdaSymbol{)}\AgdaSpace{}%
\AgdaSymbol{→}\<%
\\
%
\>[4]\AgdaFunction{interpret}\AgdaSpace{}%
\AgdaBound{c}\AgdaSpace{}%
\AgdaOperator{\AgdaFunction{≗}}\AgdaSpace{}%
\AgdaFunction{interpret}\AgdaSpace{}%
\AgdaBound{c'}\AgdaSpace{}%
\AgdaSymbol{→}\<%
\\
%
\>[4]\AgdaBound{c}\AgdaSpace{}%
\AgdaOperator{\AgdaDatatype{≡}}\AgdaSpace{}%
\AgdaBound{c'}\<%
\\
%
\>[2]\AgdaFunction{interpret-injective}\AgdaSpace{}%
\AgdaInductiveConstructor{case-constL}\AgdaSpace{}%
\AgdaInductiveConstructor{case-constL}\AgdaSpace{}%
\AgdaSymbol{\AgdaUnderscore{}}\AgdaSpace{}%
\AgdaSymbol{=}\AgdaSpace{}%
\AgdaInductiveConstructor{refl}\<%
\\
%
\>[2]\AgdaFunction{interpret-injective}\AgdaSpace{}%
\AgdaInductiveConstructor{case-constL}\AgdaSpace{}%
\AgdaInductiveConstructor{case-constR}\AgdaSpace{}%
\AgdaBound{p}\AgdaSpace{}%
\AgdaSymbol{=}\AgdaSpace{}%
\AgdaFunction{⊥-elim}\AgdaSpace{}%
\AgdaSymbol{((}\AgdaField{ℓ≠r}\AgdaSpace{}%
\AgdaBound{d}\AgdaSymbol{)}\AgdaSpace{}%
\AgdaSymbol{(}\AgdaBound{p}\AgdaSpace{}%
\AgdaSymbol{(}\AgdaField{ℓ}\AgdaSpace{}%
\AgdaBound{d}\AgdaSymbol{)))}\<%
\\
%
\>[2]\AgdaFunction{interpret-injective}\AgdaSpace{}%
\AgdaInductiveConstructor{case-constL}\AgdaSpace{}%
\AgdaInductiveConstructor{case-id}%
\>[46]\AgdaBound{p}\AgdaSpace{}%
\AgdaSymbol{=}\AgdaSpace{}%
\AgdaFunction{⊥-elim}\AgdaSpace{}%
\AgdaSymbol{((}\AgdaField{ℓ≠r}\AgdaSpace{}%
\AgdaBound{d}\AgdaSymbol{)}\AgdaSpace{}%
\AgdaSymbol{(}\AgdaBound{p}\AgdaSpace{}%
\AgdaSymbol{(}\AgdaField{r}\AgdaSpace{}%
\AgdaBound{d}\AgdaSymbol{)))}\<%
\\
%
\>[2]\AgdaFunction{interpret-injective}\AgdaSpace{}%
\AgdaInductiveConstructor{case-constL}\AgdaSpace{}%
\AgdaInductiveConstructor{case-dual}%
\>[46]\AgdaBound{p}\AgdaSpace{}%
\AgdaSymbol{=}\<%
\\
\>[2][@{}l@{\AgdaIndent{0}}]%
\>[4]\AgdaFunction{⊥-elim}\AgdaSpace{}%
\AgdaSymbol{((}\AgdaField{ℓ≠r}\AgdaSpace{}%
\AgdaBound{d}\AgdaSymbol{)}\AgdaSpace{}%
\AgdaSymbol{(}\AgdaFunction{trans}\AgdaSpace{}%
\AgdaSymbol{(}\AgdaBound{p}\AgdaSpace{}%
\AgdaSymbol{(}\AgdaField{ℓ}\AgdaSpace{}%
\AgdaBound{d}\AgdaSymbol{))}\AgdaSpace{}%
\AgdaFunction{dual-ℓ}\AgdaSymbol{))}\<%
\\
%
\>[2]\AgdaFunction{interpret-injective}\AgdaSpace{}%
\AgdaInductiveConstructor{case-constR}\AgdaSpace{}%
\AgdaInductiveConstructor{case-constL}\AgdaSpace{}%
\AgdaBound{p}\AgdaSpace{}%
\AgdaSymbol{=}\<%
\\
\>[2][@{}l@{\AgdaIndent{0}}]%
\>[4]\AgdaFunction{⊥-elim}\AgdaSpace{}%
\AgdaSymbol{((}\AgdaField{ℓ≠r}\AgdaSpace{}%
\AgdaBound{d}\AgdaSymbol{)}\AgdaSpace{}%
\AgdaSymbol{(}\AgdaFunction{sym}\AgdaSpace{}%
\AgdaSymbol{(}\AgdaBound{p}\AgdaSpace{}%
\AgdaSymbol{(}\AgdaField{ℓ}\AgdaSpace{}%
\AgdaBound{d}\AgdaSymbol{))))}\<%
\\
%
\>[2]\AgdaFunction{interpret-injective}\AgdaSpace{}%
\AgdaInductiveConstructor{case-constR}\AgdaSpace{}%
\AgdaInductiveConstructor{case-constR}\AgdaSpace{}%
\AgdaSymbol{\AgdaUnderscore{}}\AgdaSpace{}%
\AgdaSymbol{=}\AgdaSpace{}%
\AgdaInductiveConstructor{refl}\<%
\\
%
\>[2]\AgdaFunction{interpret-injective}\AgdaSpace{}%
\AgdaInductiveConstructor{case-constR}\AgdaSpace{}%
\AgdaInductiveConstructor{case-id}%
\>[46]\AgdaBound{p}\AgdaSpace{}%
\AgdaSymbol{=}\<%
\\
\>[2][@{}l@{\AgdaIndent{0}}]%
\>[4]\AgdaFunction{⊥-elim}\AgdaSpace{}%
\AgdaSymbol{((}\AgdaField{ℓ≠r}\AgdaSpace{}%
\AgdaBound{d}\AgdaSymbol{)}\AgdaSpace{}%
\AgdaSymbol{(}\AgdaFunction{sym}\AgdaSpace{}%
\AgdaSymbol{(}\AgdaBound{p}\AgdaSpace{}%
\AgdaSymbol{(}\AgdaField{ℓ}\AgdaSpace{}%
\AgdaBound{d}\AgdaSymbol{))))}\<%
\\
%
\>[2]\AgdaFunction{interpret-injective}\AgdaSpace{}%
\AgdaInductiveConstructor{case-constR}\AgdaSpace{}%
\AgdaInductiveConstructor{case-dual}%
\>[46]\AgdaBound{p}\AgdaSpace{}%
\AgdaSymbol{=}\<%
\\
\>[2][@{}l@{\AgdaIndent{0}}]%
\>[4]\AgdaFunction{⊥-elim}\AgdaSpace{}%
\AgdaSymbol{((}\AgdaField{ℓ≠r}\AgdaSpace{}%
\AgdaBound{d}\AgdaSymbol{)}\AgdaSpace{}%
\AgdaSymbol{(}\AgdaFunction{sym}\AgdaSpace{}%
\AgdaSymbol{(}\AgdaFunction{trans}\AgdaSpace{}%
\AgdaSymbol{(}\AgdaBound{p}\AgdaSpace{}%
\AgdaSymbol{(}\AgdaField{r}\AgdaSpace{}%
\AgdaBound{d}\AgdaSymbol{))}\AgdaSpace{}%
\AgdaFunction{dual-r}\AgdaSymbol{)))}\<%
\\
%
\>[2]\AgdaFunction{interpret-injective}\AgdaSpace{}%
\AgdaInductiveConstructor{case-id}%
\>[34]\AgdaInductiveConstructor{case-constL}\AgdaSpace{}%
\AgdaBound{p}\AgdaSpace{}%
\AgdaSymbol{=}\<%
\\
\>[2][@{}l@{\AgdaIndent{0}}]%
\>[4]\AgdaFunction{⊥-elim}\AgdaSpace{}%
\AgdaSymbol{((}\AgdaField{ℓ≠r}\AgdaSpace{}%
\AgdaBound{d}\AgdaSymbol{)}\AgdaSpace{}%
\AgdaSymbol{(}\AgdaFunction{sym}\AgdaSpace{}%
\AgdaSymbol{(}\AgdaBound{p}\AgdaSpace{}%
\AgdaSymbol{(}\AgdaField{r}\AgdaSpace{}%
\AgdaBound{d}\AgdaSymbol{))))}\<%
\\
%
\>[2]\AgdaFunction{interpret-injective}\AgdaSpace{}%
\AgdaInductiveConstructor{case-id}%
\>[34]\AgdaInductiveConstructor{case-constR}\AgdaSpace{}%
\AgdaBound{p}\AgdaSpace{}%
\AgdaSymbol{=}\AgdaSpace{}%
\AgdaFunction{⊥-elim}\AgdaSpace{}%
\AgdaSymbol{((}\AgdaField{ℓ≠r}\AgdaSpace{}%
\AgdaBound{d}\AgdaSymbol{)}\AgdaSpace{}%
\AgdaSymbol{(}\AgdaBound{p}\AgdaSpace{}%
\AgdaSymbol{(}\AgdaField{ℓ}\AgdaSpace{}%
\AgdaBound{d}\AgdaSymbol{)))}\<%
\\
%
\>[2]\AgdaFunction{interpret-injective}\AgdaSpace{}%
\AgdaInductiveConstructor{case-id}%
\>[34]\AgdaInductiveConstructor{case-id}%
\>[46]\AgdaSymbol{\AgdaUnderscore{}}\AgdaSpace{}%
\AgdaSymbol{=}\AgdaSpace{}%
\AgdaInductiveConstructor{refl}\<%
\\
%
\>[2]\AgdaFunction{interpret-injective}\AgdaSpace{}%
\AgdaInductiveConstructor{case-id}%
\>[34]\AgdaInductiveConstructor{case-dual}%
\>[46]\AgdaBound{p}\AgdaSpace{}%
\AgdaSymbol{=}\<%
\\
\>[2][@{}l@{\AgdaIndent{0}}]%
\>[4]\AgdaFunction{⊥-elim}\AgdaSpace{}%
\AgdaSymbol{((}\AgdaField{ℓ≠r}\AgdaSpace{}%
\AgdaBound{d}\AgdaSymbol{)}\AgdaSpace{}%
\AgdaSymbol{(}\AgdaFunction{trans}\AgdaSpace{}%
\AgdaSymbol{(}\AgdaBound{p}\AgdaSpace{}%
\AgdaSymbol{(}\AgdaField{ℓ}\AgdaSpace{}%
\AgdaBound{d}\AgdaSymbol{))}\AgdaSpace{}%
\AgdaFunction{dual-ℓ}\AgdaSymbol{))}\<%
\\
%
\>[2]\AgdaFunction{interpret-injective}\AgdaSpace{}%
\AgdaInductiveConstructor{case-dual}%
\>[34]\AgdaInductiveConstructor{case-constL}\AgdaSpace{}%
\AgdaBound{p}\AgdaSpace{}%
\AgdaSymbol{=}\<%
\\
\>[2][@{}l@{\AgdaIndent{0}}]%
\>[4]\AgdaFunction{⊥-elim}\AgdaSpace{}%
\AgdaSymbol{((}\AgdaField{ℓ≠r}\AgdaSpace{}%
\AgdaBound{d}\AgdaSymbol{)}\AgdaSpace{}%
\AgdaSymbol{(}\AgdaFunction{sym}\AgdaSpace{}%
\AgdaSymbol{(}\AgdaFunction{trans}\AgdaSpace{}%
\AgdaSymbol{(}\AgdaFunction{sym}\AgdaSpace{}%
\AgdaSymbol{(}\AgdaFunction{dual-ℓ}\AgdaSymbol{))}\AgdaSpace{}%
\AgdaSymbol{(}\AgdaBound{p}\AgdaSpace{}%
\AgdaSymbol{(}\AgdaField{ℓ}\AgdaSpace{}%
\AgdaBound{d}\AgdaSymbol{)))))}\<%
\\
%
\>[2]\AgdaFunction{interpret-injective}\AgdaSpace{}%
\AgdaInductiveConstructor{case-dual}%
\>[34]\AgdaInductiveConstructor{case-constR}\AgdaSpace{}%
\AgdaBound{p}\AgdaSpace{}%
\AgdaSymbol{=}\<%
\\
\>[2][@{}l@{\AgdaIndent{0}}]%
\>[4]\AgdaFunction{⊥-elim}\AgdaSpace{}%
\AgdaSymbol{((}\AgdaField{ℓ≠r}\AgdaSpace{}%
\AgdaBound{d}\AgdaSymbol{)}\AgdaSpace{}%
\AgdaSymbol{(}\AgdaFunction{trans}\AgdaSpace{}%
\AgdaSymbol{(}\AgdaFunction{sym}\AgdaSpace{}%
\AgdaSymbol{(}\AgdaFunction{dual-r}\AgdaSymbol{))}\AgdaSpace{}%
\AgdaSymbol{(}\AgdaBound{p}\AgdaSpace{}%
\AgdaSymbol{(}\AgdaField{r}\AgdaSpace{}%
\AgdaBound{d}\AgdaSymbol{))))}\<%
\\
%
\>[2]\AgdaFunction{interpret-injective}\AgdaSpace{}%
\AgdaInductiveConstructor{case-dual}%
\>[34]\AgdaInductiveConstructor{case-id}%
\>[46]\AgdaBound{p}\AgdaSpace{}%
\AgdaSymbol{=}\<%
\\
\>[2][@{}l@{\AgdaIndent{0}}]%
\>[4]\AgdaFunction{⊥-elim}\AgdaSpace{}%
\AgdaSymbol{((}\AgdaField{ℓ≠r}\AgdaSpace{}%
\AgdaBound{d}\AgdaSymbol{)}\AgdaSpace{}%
\AgdaSymbol{(}\AgdaFunction{sym}\AgdaSpace{}%
\AgdaSymbol{(}\AgdaFunction{trans}\AgdaSpace{}%
\AgdaSymbol{(}\AgdaFunction{sym}\AgdaSpace{}%
\AgdaSymbol{(}\AgdaFunction{dual-ℓ}\AgdaSymbol{))}\AgdaSpace{}%
\AgdaSymbol{(}\AgdaBound{p}\AgdaSpace{}%
\AgdaSymbol{(}\AgdaField{ℓ}\AgdaSpace{}%
\AgdaBound{d}\AgdaSymbol{)))))}\<%
\\
%
\>[2]\AgdaFunction{interpret-injective}\AgdaSpace{}%
\AgdaInductiveConstructor{case-dual}%
\>[34]\AgdaInductiveConstructor{case-dual}%
\>[46]\AgdaSymbol{\AgdaUnderscore{}}\AgdaSpace{}%
\AgdaSymbol{=}\AgdaSpace{}%
\AgdaInductiveConstructor{refl}\<%
\end{code}

\begin{code}%
%
\>[2]\AgdaFunction{classify-unique}\AgdaSpace{}%
\AgdaSymbol{:}\AgdaSpace{}%
\AgdaSymbol{(}\AgdaBound{f}\AgdaSpace{}%
\AgdaSymbol{:}\AgdaSpace{}%
\AgdaFunction{Endo}\AgdaSymbol{)}\AgdaSpace{}%
\AgdaSymbol{→}\AgdaSpace{}%
\AgdaSymbol{(}\AgdaBound{c}\AgdaSpace{}%
\AgdaSymbol{:}\AgdaSpace{}%
\AgdaDatatype{EndoCase}\AgdaSymbol{)}\AgdaSpace{}%
\AgdaSymbol{→}\AgdaSpace{}%
\AgdaFunction{interpret}\AgdaSpace{}%
\AgdaBound{c}\AgdaSpace{}%
\AgdaOperator{\AgdaFunction{≗}}\AgdaSpace{}%
\AgdaBound{f}\AgdaSpace{}%
\AgdaSymbol{→}\AgdaSpace{}%
\AgdaBound{c}\AgdaSpace{}%
\AgdaOperator{\AgdaDatatype{≡}}\AgdaSpace{}%
\AgdaFunction{classify}\AgdaSpace{}%
\AgdaBound{f}\<%
\\
%
\>[2]\AgdaFunction{classify-unique}\AgdaSpace{}%
\AgdaBound{f}\AgdaSpace{}%
\AgdaBound{c}\AgdaSpace{}%
\AgdaBound{c≗f}\AgdaSpace{}%
\AgdaSymbol{=}\<%
\\
\>[2][@{}l@{\AgdaIndent{0}}]%
\>[4]\AgdaFunction{interpret-injective}\AgdaSpace{}%
\AgdaBound{c}\AgdaSpace{}%
\AgdaSymbol{(}\AgdaFunction{classify}\AgdaSpace{}%
\AgdaBound{f}\AgdaSymbol{)}\AgdaSpace{}%
\AgdaSymbol{(}\AgdaFunction{≗-trans}\AgdaSpace{}%
\AgdaBound{c≗f}\AgdaSpace{}%
\AgdaSymbol{(}\AgdaFunction{≗-sym}\AgdaSpace{}%
\AgdaSymbol{(}\AgdaFunction{classify-sound}\AgdaSpace{}%
\AgdaBound{f}\AgdaSymbol{)))}\<%
\\
%
\\[\AgdaEmptyExtraSkip]%
%
\>[2]\AgdaFunction{classify-interpret}\AgdaSpace{}%
\AgdaSymbol{:}\AgdaSpace{}%
\AgdaSymbol{(}\AgdaBound{c}\AgdaSpace{}%
\AgdaSymbol{:}\AgdaSpace{}%
\AgdaDatatype{EndoCase}\AgdaSymbol{)}\AgdaSpace{}%
\AgdaSymbol{→}\AgdaSpace{}%
\AgdaFunction{classify}\AgdaSpace{}%
\AgdaSymbol{(}\AgdaFunction{interpret}\AgdaSpace{}%
\AgdaBound{c}\AgdaSymbol{)}\AgdaSpace{}%
\AgdaOperator{\AgdaDatatype{≡}}\AgdaSpace{}%
\AgdaBound{c}\<%
\\
%
\>[2]\AgdaFunction{classify-interpret}\AgdaSpace{}%
\AgdaBound{c}\AgdaSpace{}%
\AgdaSymbol{=}\AgdaSpace{}%
\AgdaFunction{sym}\AgdaSpace{}%
\AgdaSymbol{(}\AgdaFunction{classify-unique}\AgdaSpace{}%
\AgdaSymbol{(}\AgdaFunction{interpret}\AgdaSpace{}%
\AgdaBound{c}\AgdaSymbol{)}\AgdaSpace{}%
\AgdaBound{c}\AgdaSpace{}%
\AgdaFunction{≗-refl}\AgdaSymbol{)}\<%
\end{code}

Together this says: the four labels are pairwise distinct, and every
self-map falls uniquely under one of them.

\begin{code}%
\>[0]\AgdaFunction{EndoCase-distinct}\AgdaSpace{}%
\AgdaSymbol{:}\<%
\\
\>[0][@{}l@{\AgdaIndent{0}}]%
\>[2]\AgdaSymbol{(}\AgdaInductiveConstructor{case-constL}\AgdaSpace{}%
\AgdaOperator{\AgdaDatatype{≡}}\AgdaSpace{}%
\AgdaInductiveConstructor{case-constR}\AgdaSpace{}%
\AgdaSymbol{→}\AgdaSpace{}%
\AgdaDatatype{⊥}\AgdaSymbol{)}\AgdaSpace{}%
\AgdaOperator{\AgdaRecord{×}}\<%
\\
%
\>[2]\AgdaSymbol{(}\AgdaInductiveConstructor{case-constL}\AgdaSpace{}%
\AgdaOperator{\AgdaDatatype{≡}}\AgdaSpace{}%
\AgdaInductiveConstructor{case-id}%
\>[29]\AgdaSymbol{→}\AgdaSpace{}%
\AgdaDatatype{⊥}\AgdaSymbol{)}\AgdaSpace{}%
\AgdaOperator{\AgdaRecord{×}}\<%
\\
%
\>[2]\AgdaSymbol{(}\AgdaInductiveConstructor{case-constL}\AgdaSpace{}%
\AgdaOperator{\AgdaDatatype{≡}}\AgdaSpace{}%
\AgdaInductiveConstructor{case-dual}%
\>[29]\AgdaSymbol{→}\AgdaSpace{}%
\AgdaDatatype{⊥}\AgdaSymbol{)}\AgdaSpace{}%
\AgdaOperator{\AgdaRecord{×}}\<%
\\
%
\>[2]\AgdaSymbol{(}\AgdaInductiveConstructor{case-constR}\AgdaSpace{}%
\AgdaOperator{\AgdaDatatype{≡}}\AgdaSpace{}%
\AgdaInductiveConstructor{case-id}%
\>[29]\AgdaSymbol{→}\AgdaSpace{}%
\AgdaDatatype{⊥}\AgdaSymbol{)}\AgdaSpace{}%
\AgdaOperator{\AgdaRecord{×}}\<%
\\
%
\>[2]\AgdaSymbol{(}\AgdaInductiveConstructor{case-constR}\AgdaSpace{}%
\AgdaOperator{\AgdaDatatype{≡}}\AgdaSpace{}%
\AgdaInductiveConstructor{case-dual}%
\>[29]\AgdaSymbol{→}\AgdaSpace{}%
\AgdaDatatype{⊥}\AgdaSymbol{)}\AgdaSpace{}%
\AgdaOperator{\AgdaRecord{×}}\<%
\\
%
\>[2]\AgdaSymbol{(}\AgdaInductiveConstructor{case-id}%
\>[15]\AgdaOperator{\AgdaDatatype{≡}}\AgdaSpace{}%
\AgdaInductiveConstructor{case-dual}%
\>[29]\AgdaSymbol{→}\AgdaSpace{}%
\AgdaDatatype{⊥}\AgdaSymbol{)}\<%
\\
\>[0]\AgdaFunction{EndoCase-distinct}\AgdaSpace{}%
\AgdaSymbol{=}\<%
\\
\>[0][@{}l@{\AgdaIndent{0}}]%
\>[2]\AgdaSymbol{(λ}\AgdaSpace{}%
\AgdaSymbol{())}\AgdaSpace{}%
\AgdaOperator{\AgdaInductiveConstructor{,}}\<%
\\
%
\>[2]\AgdaSymbol{((λ}\AgdaSpace{}%
\AgdaSymbol{())}\AgdaSpace{}%
\AgdaOperator{\AgdaInductiveConstructor{,}}\<%
\\
\>[2][@{}l@{\AgdaIndent{0}}]%
\>[3]\AgdaSymbol{((λ}\AgdaSpace{}%
\AgdaSymbol{())}\AgdaSpace{}%
\AgdaOperator{\AgdaInductiveConstructor{,}}\<%
\\
\>[3][@{}l@{\AgdaIndent{0}}]%
\>[4]\AgdaSymbol{((λ}\AgdaSpace{}%
\AgdaSymbol{())}\AgdaSpace{}%
\AgdaOperator{\AgdaInductiveConstructor{,}}\<%
\\
\>[4][@{}l@{\AgdaIndent{0}}]%
\>[5]\AgdaSymbol{((λ}\AgdaSpace{}%
\AgdaSymbol{())}\AgdaSpace{}%
\AgdaOperator{\AgdaInductiveConstructor{,}}\<%
\\
\>[5][@{}l@{\AgdaIndent{0}}]%
\>[6]\AgdaSymbol{(λ}\AgdaSpace{}%
\AgdaSymbol{())))))}\<%
\end{code}

Gathered into two top-level statements, the result reads cleanly: every
endomorphism of a distinction is one of the four cases (soundly), and is
so in exactly one way (uniquely).

\begin{code}%
\>[0]\AgdaFunction{k4-classification-sound}\AgdaSpace{}%
\AgdaSymbol{:}\<%
\\
\>[0][@{}l@{\AgdaIndent{0}}]%
\>[2]\AgdaSymbol{(}\AgdaBound{d}\AgdaSpace{}%
\AgdaSymbol{:}\AgdaSpace{}%
\AgdaRecord{Distinction}\AgdaSymbol{)}\AgdaSpace{}%
\AgdaSymbol{→}\<%
\\
%
\>[2]\AgdaSymbol{(}\AgdaBound{f}\AgdaSpace{}%
\AgdaSymbol{:}\AgdaSpace{}%
\AgdaField{S}\AgdaSpace{}%
\AgdaBound{d}\AgdaSpace{}%
\AgdaSymbol{→}\AgdaSpace{}%
\AgdaField{S}\AgdaSpace{}%
\AgdaBound{d}\AgdaSymbol{)}\AgdaSpace{}%
\AgdaSymbol{→}\<%
\\
%
\>[2]\AgdaRecord{Σ}\AgdaSpace{}%
\AgdaDatatype{EndoCase}\AgdaSpace{}%
\AgdaSymbol{(λ}\AgdaSpace{}%
\AgdaBound{c}\AgdaSpace{}%
\AgdaSymbol{→}\AgdaSpace{}%
\AgdaFunction{K₄.interpret}\AgdaSpace{}%
\AgdaBound{d}\AgdaSpace{}%
\AgdaBound{c}\AgdaSpace{}%
\AgdaOperator{\AgdaFunction{≗}}\AgdaSpace{}%
\AgdaBound{f}\AgdaSymbol{)}\<%
\\
\>[0]\AgdaFunction{k4-classification-sound}\AgdaSpace{}%
\AgdaBound{d}\AgdaSpace{}%
\AgdaBound{f}\AgdaSpace{}%
\AgdaSymbol{=}\AgdaSpace{}%
\AgdaFunction{K₄.classify}\AgdaSpace{}%
\AgdaBound{d}\AgdaSpace{}%
\AgdaBound{f}\AgdaSpace{}%
\AgdaOperator{\AgdaInductiveConstructor{,}}\AgdaSpace{}%
\AgdaFunction{K₄.classify-sound}\AgdaSpace{}%
\AgdaBound{d}\AgdaSpace{}%
\AgdaBound{f}\<%
\\
%
\\[\AgdaEmptyExtraSkip]%
\>[0]\AgdaFunction{k4-classification-unique}\AgdaSpace{}%
\AgdaSymbol{:}\<%
\\
\>[0][@{}l@{\AgdaIndent{0}}]%
\>[2]\AgdaSymbol{(}\AgdaBound{d}\AgdaSpace{}%
\AgdaSymbol{:}\AgdaSpace{}%
\AgdaRecord{Distinction}\AgdaSymbol{)}\AgdaSpace{}%
\AgdaSymbol{→}\<%
\\
%
\>[2]\AgdaSymbol{(}\AgdaBound{f}\AgdaSpace{}%
\AgdaSymbol{:}\AgdaSpace{}%
\AgdaField{S}\AgdaSpace{}%
\AgdaBound{d}\AgdaSpace{}%
\AgdaSymbol{→}\AgdaSpace{}%
\AgdaField{S}\AgdaSpace{}%
\AgdaBound{d}\AgdaSymbol{)}\AgdaSpace{}%
\AgdaSymbol{→}\<%
\\
%
\>[2]\AgdaSymbol{(}\AgdaBound{c₁}\AgdaSpace{}%
\AgdaBound{c₂}\AgdaSpace{}%
\AgdaSymbol{:}\AgdaSpace{}%
\AgdaDatatype{EndoCase}\AgdaSymbol{)}\AgdaSpace{}%
\AgdaSymbol{→}\<%
\\
%
\>[2]\AgdaFunction{K₄.interpret}\AgdaSpace{}%
\AgdaBound{d}\AgdaSpace{}%
\AgdaBound{c₁}\AgdaSpace{}%
\AgdaOperator{\AgdaFunction{≗}}\AgdaSpace{}%
\AgdaBound{f}\AgdaSpace{}%
\AgdaSymbol{→}\<%
\\
%
\>[2]\AgdaFunction{K₄.interpret}\AgdaSpace{}%
\AgdaBound{d}\AgdaSpace{}%
\AgdaBound{c₂}\AgdaSpace{}%
\AgdaOperator{\AgdaFunction{≗}}\AgdaSpace{}%
\AgdaBound{f}\AgdaSpace{}%
\AgdaSymbol{→}\<%
\\
%
\>[2]\AgdaBound{c₁}\AgdaSpace{}%
\AgdaOperator{\AgdaDatatype{≡}}\AgdaSpace{}%
\AgdaBound{c₂}\<%
\\
\>[0]\AgdaFunction{k4-classification-unique}\AgdaSpace{}%
\AgdaBound{d}\AgdaSpace{}%
\AgdaBound{f}\AgdaSpace{}%
\AgdaBound{c₁}\AgdaSpace{}%
\AgdaBound{c₂}\AgdaSpace{}%
\AgdaBound{p₁}\AgdaSpace{}%
\AgdaBound{p₂}\AgdaSpace{}%
\AgdaSymbol{=}\<%
\\
\>[0][@{}l@{\AgdaIndent{0}}]%
\>[2]\AgdaFunction{K₄.interpret-injective}\AgdaSpace{}%
\AgdaBound{d}\AgdaSpace{}%
\AgdaBound{c₁}\AgdaSpace{}%
\AgdaBound{c₂}\AgdaSpace{}%
\AgdaSymbol{(}\AgdaFunction{K₄.≗-trans}\AgdaSpace{}%
\AgdaBound{d}\AgdaSpace{}%
\AgdaBound{p₁}\AgdaSpace{}%
\AgdaSymbol{(}\AgdaFunction{K₄.≗-sym}\AgdaSpace{}%
\AgdaBound{d}\AgdaSpace{}%
\AgdaBound{p₂}\AgdaSymbol{))}\<%
\end{code}

A quick algebraic over-reading is tempting. But what has been shown is only
this: the four maps are closed under composition, and they do not form a
group because the two constants are not invertible. What is established is
the number of cases and their unique classification up to pointwise
equality. Any richer algebraic reading --- and likewise the later reading
of this four-case structure as a $K_4$ closure --- is an additional step.

The next answer has therefore been reached. From a separated and covered
binary distinction follows a unique four-case classification of its
self-maps. The next question no longer concerns the structure itself, but
its status: what has thereby been shown, and what not?

% ─────────────────────────────────────────────────────────────────────
\chapter{What Has Been Shown}
\label{ch:reading-the-threshold}
\label{ch:toward-form}
\label{ch:conclusion}
% ─────────────────────────────────────────────────────────────────────

One sentence must now be said as sharply as possible. This book has not
proved the world. It has shown that a stable distinction, once written
exactly under the stated constraints, already carries a nontrivial closed
structure.

What has been shown is a small but precise stock: a carrier with two
positions, a witness that they are apart, an exhaustion principle with no
third case, and exactly four self-maps up to pointwise equality. Under
\AgdaFlag{safe} and \AgdaFlag{without-K}, these statements are theorems in
the strict sense: if any one of them failed, the file would not compile.

At this point the levels of the book can be ordered cleanly. Philosophical
orientation was the conjecture that origin, determination, and information
all point toward distinction. The working assumption was that successful
distinction must satisfy minimal conditions. The formal result is the stock
the file actually carries. It is on that level that this chapter moves.

What has not been shown is that the world can be derived from nothing, or
that every determination necessarily presupposes distinction as a theorem
of Agda. Those claims motivate the inquiry; they are not its checked
result. Nor does the file show that no other formal carrier could realize
the same content. It shows only that this content can be constructed here
without additional assumptions.

The result is nonetheless interesting. Once two positions are granted in
full view and kept apart, that minimal grant already carries a closed
shape. This is the first point at which the conditions of stable
distinction become exact and machine-checkable. The step from
philosophical question to formal structure remains small. That is exactly
why it is robust.

From here one can also state the boundaries cleanly. The compiler judges
the formal kernel. Ordinary mathematics judges the derivations built from
that kernel. Experiment judges any later reading of the kernel as physics.
The separation between those courts is not caution for its own sake. It is
a condition of honesty.

At the level of this book, the status is clear. That a stable distinction
has exactly four self-maps up to pointwise equality is theorem. Further
constructions remain mathematics so long as they remain explicit
constructions. Only when a later stage gives such structures physical names
does theorem give way to interpretation, and only confrontation with the
world can decide whether an interpretation survives.

The question of origin therefore remains unanswered. What has shifted is
narrower and more precise: its first minimal step --- the holding-apart of
one position from another --- is no longer mere intuition. It is present as
checked construction, bounded carrier, and four-case self-map structure.

% ─────────────────────────────────────────────────────────────────────
\chapter{The Reach of Distinction}
\label{ch:outlook}
\label{ch:reach}
% ─────────────────────────────────────────────────────────────────────

This companion is not the theory as a whole. It is not the physics. It is
not the unification. It marks the first point at which the conditions of
stable distinction can be formulated completely and checked mechanically.

From here the larger investigation begins rather than ends.
\textit{The First Distinction} establishes the datum. \textit{Void} asks
what follows if nothing further is granted. It does not ask which external
mathematics can be borrowed, but how much mathematics can be reconstructed
from the same discipline. Only after that register exists does
\textit{Form} begin. There the question is no longer what follows formally,
but whether structures already present in the ledger can be read as
structures of physical description.

Above both stands the wider question that motivates the whole project:
whether logic, mathematics, and physical description may be different
appearances of a common origin.

\begin{center}
\begin{tikzpicture}[node distance=8mm]
  \node[stage, minimum width=44mm] (fd) {The First\\Distinction};
  \node[stage, minimum width=44mm, below=of fd] (void) {Void};
  \node[stage, minimum width=44mm, below=of void] (math) {Mathematical\\Reconstruction};
  \node[stage, minimum width=44mm, below=of math] (phys) {Physical\\Reconstruction};
  \node[stage, minimum width=44mm, below=of phys] (common) {Common-Origin\\Question};

  \draw[routearrow] (fd) -- (void);
  \draw[routearrow] (void) -- (math);
  \draw[routearrow] (math) -- (phys);
  \draw[routearrow] (phys) -- (common);
\end{tikzpicture}
\end{center}

Void/Form investigates whether logic, mathematics, and physical description
can be understood as different appearances of a common structure of
distinction. \textit{The First Distinction} supplies the first formally
secured step of that investigation.

How far can this go? This book does not answer that question. It only makes
the question exact enough to be pursued honestly. The chain may break
early, and then its limit will have been identified. Or it may carry far,
and then its reach will have been measured rather than proclaimed.

Where a later physical reading is proposed, the formal system does not
decide it. There the ledger meets experiment, and science begins in the
strict sense.

This book therefore ends where the larger inquiry truly begins: with the
first distinction in hand, and with the reach of that distinction still
open.

\end{document}
